\documentclass[11pt]{article}
\usepackage[letterpaper,margin=1in]{geometry}
\usepackage[T1]{fontenc}
\usepackage{lmodern}
\usepackage{amsmath,amssymb,amsthm,mathtools}
\usepackage{mathrsfs}
\usepackage{microtype}
\usepackage{xcolor}
\usepackage{hyperref}
\definecolor{linkblue}{RGB}{28,65,105}
\hypersetup{colorlinks=true,linkcolor=linkblue,citecolor=linkblue,urlcolor=linkblue,
 pdftitle={Conservative Theta Flows and Collision Amplification},
 pdfauthor={Drafted for Edward Baker}}
\numberwithin{equation}{section}
\newtheorem{theorem}{Theorem}[section]
\newtheorem{proposition}[theorem]{Proposition}
\newtheorem{lemma}[theorem]{Lemma}
\newcommand{\dd}{\,\mathrm d}
\newcommand{\R}{\mathbb R}
\setlength{\emergencystretch}{2em}
\title{Conservative Theta Flows and Collision Amplification}
\author{Drafted for Edward Baker}
\date{29 October 2026\\[3pt]\small Working manuscript}

\begin{document}
\maketitle

\begin{abstract}
We study conservative flows of the complete theta observation and its
actual frequency derivative, allowing their densities to redistribute
across radii. Ordinary spatial heat preserves the entire integral and
gives an explicit central decay diagnostic for simultaneous cancellation.
An energy-normalized conservative feedback fixes the joint spatial
$L^2$ energy of the observation and its marked first derivative. A
collision forces gradient, drift-response, and amplitude divergence
before an explicit finite auxiliary deadline; global regularity is
equivalent to pointwise noncollision in the stated class. A sharp
positive-carrier comparison
gives quantitative regularity only from an independently established
first-jet margin. The complete canonical source relation instead
acquires a nonzero preparation residual and second-order commutator
forcing. Its unit-sphere cancellation, inner-ball energy, and admissible
origin domain must be retained. A weighted energy coupling the gradient,
residual and actual second spectral derivative absorbs the commutator's
second spatial derivatives into diffusion. It yields integrated
dissipation and a local existence bound without a third spectral jet.
An exact radial drift square further lowers the gradient-production
energy allowance from $30E_0$ to $21E_0/4$ and gives a longer local
interval, including comparison with the old bound retaining diffusion.
An all-radius cutoff estimate also makes the compact weighted
coefficient gain strict, with gap at least $41/24$.
Coupling the complete source coefficient to the actual residual lowers
that compact coefficient by a further $49/600$, with all jet and inner
costs paid. This source-coefficient gain still admits the exact colliding
control. These bounds retain no separating arithmetic sign and do not control
the collision deadline. The exact colliding control
satisfies the generic heat, feedback and weighted-estimate premises. No new theta
collision exclusion or de Bruijn--Newman bound is obtained.
\end{abstract}

\section{Purpose and conserved observation}

The invariant is the complete theta observation $H(t,a)$, together with
its actual frequency derivative $H_a(t,a)$. The collision question is
whether these two quantities can vanish simultaneously at a physical
mark $(T,A)$. Our strategy is to preserve their integrals, magnify their
simultaneous cancellation by a conservative evolution, and seek an
independent estimate from the unchanged complete theta source that
controls the magnifier through its collision deadline.

The admissible search permits redistribution across radii. The theta
heat and feedback solutions considered here remain radially symmetric,
but their radial profiles change with auxiliary time. Radial symmetry
does not preserve the radial integrand. The density, including its
frequency dependence, is what evolves on all of $\R^3$; its initially
empty unit ball belongs to the evolution.

After the source setup, Section~\ref{sec:conservative} constructs the
whole-integral heat flow, states its exact colliding control, and derives
the central decay diagnostic \cite{ConservativeNote}.
Section~\ref{sec:feedback} fixes the energy-normalized feedback and its
finite collision deadline \cite{FeedbackNote}.
Sections~\ref{sec:carrier} and \ref{sec:sourcerelation} test,
respectively, positive-carrier comparison and the canonical source
relation with its residual \cite{CarrierNote,SourceRelationNote}.
Section~\ref{sec:weightedclosure} proves finite weighted residual closure
\cite{WeightedClosureNote}. Section~\ref{sec:radialgain} sharpens its
gradient-production allowance by an exact radial square
\cite{RadialDriftNote}. Section~\ref{sec:sourcecoupling} retains the
source--residual product for a further compact gain \cite{SourceCouplingNote}.
Section~\ref{sec:remaining} identifies the
remaining improvement and collision-deadline obligations.
The classical heat and semigroup inputs support these constructions;
the missing estimate must use canonical information beyond their
generic premises.

We use $r=e^{2u}$, with $u=\frac12\log r$ on the source exterior,
$a$ for the physical frequency, $t$ for Newman time, and $s\ge0$ for
auxiliary time. Newman time satisfies $H_t=-H_{aa}$. In the feedback
construction the mark $p=(T,A)$ is fixed while the running parameters
$(t,a)$ vary. A spectral derivative differentiates $a$ with $p$ held
fixed; changing $A$ also changes the feedback path. Neither physical
parameter evolves with $s$.
Unless a wider scope is stated, uniform assertions concern compact sets
\begin{equation}\label{eq:domain}
0\le t\le T_+<\infty,\qquad
0<a_-\le a\le a_+<\infty.
\end{equation}
The source-function, smoothing, and operator domains, and the control's
convergence domain, are specified where they are used.

The remaining estimate can be stated before the construction. For the
marked pair $W=(U,V)$ of Section~\ref{sec:feedback}, write
$G(s)=\|\nabla W_s\|_2^2$. The target is an independent bound for
$G$, or its feedback response $\alpha$, that keeps the solution regular
through the explicit collision deadline $S_*$ in
Theorem~\ref{thm:feedbackdeadline}. Its exact gradient balance is
\eqref{eq:feedbackgradientexact}. A bounded improvement in gradient or
residual control is an intermediate result if its consequence for that
balance is quantified. Section~\ref{sec:weightedclosure} establishes such
a weighted residual estimate and a generic radial gradient improvement,
while deadline transfer remains open.
A collision-excluding bound must use the complete theta source
without assuming the nonzero physical first jet that it is meant to
establish. Sections~\ref{sec:carrier}--\ref{sec:remaining} retain the
tested estimates and the remaining gap.

\section{The canonical theta source and observation}

Define
\begin{align}
\Theta(r)&=1+2\sum_{n\ge1}e^{-\pi n^2r^2},
&\Delta_3&=\partial_r^2+\frac2r\partial_r,\label{eq:theta}\\
Q(r)&=\Delta_3\Theta(r)
=4\pi\sum_{n\ge1}n^2(2\pi n^2r^2-3)e^{-\pi n^2r^2},\label{eq:Q}\\
\Phi(u)&=\frac14r^{5/2}Q(r),
&Q_t(r)&=e^{t\log^2r/4}Q(r),\qquad r=e^{2u}.\label{eq:fullthetaflow}
\end{align}
For $r\ge1$, every summand in \eqref{eq:Q} is positive. The complete
physical observation is
\begin{equation}\label{eq:H}
H(t,a)=\int_0^\infty e^{tu^2}\Phi(u)\cos(au)\dd u
=\frac18\int_1^\infty r^{3/2}Q_t(r)
             \cos\!\left(\frac a2\log r\right)\dd r.
\end{equation}
The second formula already represents the full first integral. An
inversion-related interior integral is not to be added again.
The transform companion derives the related classical theta--Mellin
formula in this normalization \cite{ThetaCompanion}; its potential and
cylinder fields are background representations, while the density
specified in Section~\ref{sec:conservative} is the field evolved here.

\begin{lemma}[Tails and parameter differentiation]\label{lem:tails}
The canonical source and the finite spatial and parameter derivatives
used below have integrable majorants on \eqref{eq:domain}. In the
coordinate $u\ge0$, they decay faster than every fixed exponential,
uniformly on compact parameter sets.
\end{lemma}
\begin{proof}
For an integer $p\ge0$ and $r\ge1$,
\[
\sum_{n\ge1}n^p e^{-\pi n^2r^2}
\le C_p e^{-\pi r^2},\qquad
C_p=\sum_{j\ge0}(1+j)^p e^{-3\pi j}<\infty,
\]
because $(1+j)^2\ge1+3j$. Finite differentiation adds polynomial
powers of $n,r$ and $\log r$. These are absorbed by a slight decrease
of the positive Gaussian exponent. For the time multiplier, the explicit
inequalities
\[
\log^2r\le\frac{4r}{e^2},\qquad
\frac{T_+r}{e^2}\le\frac\pi2r^2+\frac{T_+^2}{2\pi e^4}
\]
give a uniform bound by a constant times $e^{-\pi r^2/2}$.
Since $r=e^{2u}$, this also proves the asserted exponential domination in
$u$. It justifies the differentiations and outer boundary limits below.
\end{proof}

\section{Conserving the entire integral}\label{sec:conservative}

\subsection{A local conservation equation}

The chosen invariant is represented by a density $j_s$ with
\begin{equation}\label{eq:integralconservation}
\int_\Omega j_s(t,a;X)\dd X=H(t,a).
\end{equation}
Its radial average may change. A local equation is
\begin{equation}\label{eq:continuity}
\partial_sj_s+\nabla\cdot\mathbf J_s=0.
\end{equation}
The divergence theorem preserves \eqref{eq:integralconservation} when
the complete boundary flux, including any outer limit, is zero.
A subregion instead gains or loses precisely its boundary flux.

To distinguish this evolution from diffusion of $Q_t$ alone, define
the spherical average of an exterior source by
\[
 \mathcal P F_s(r)=\frac1{4\pi}\int_{S^2}F_s(r\omega)\dd\omega,
 \qquad r\ge1.
\]
If a frequency-independent $F_s$ is tested against the same fixed
kernel $r^{-1/2}\cos(a\log r/2)$ at every real $a$, invariance of those
exterior integrals forces $\mathcal P F_s=Q_t$, provided
$r^{3/2}(\mathcal P F_s-Q_t)\in L^1([1,\infty))$.
Indeed $r=e^{2u}$ reduces the difference to an integrable cosine
transform; Gaussian convolution of its even extension proves uniqueness.
Evolving the full frequency-dependent density, as below, permits radial
redistribution without that fixed-source constraint. Moving the testing
kernel together with the source is another possibility. No division by
the cosine at its zeros is appropriate.

For example, transporting a density and phase by
\[
\partial_s\rho+\nabla\cdot(\rho v)=0,\qquad
\partial_s\ell+v\cdot\nabla\ell=0
\]
makes $\rho_s\cos(a\ell_s)$ conservative wherever the smooth flow and
its boundary conditions are defined. More simply, choose the diffusive
flux $\mathbf J_s=-\varepsilon\nabla j_s$.

\subsection{Heat flow of the full theta integrand}

On $\mathbb R^3$, start from
\begin{equation}\label{eq:densityinitial}
j_0(t,a;X)=
 \begin{cases}
 (32\pi)^{-1}r^{-1/2}Q_t(r)\cos(a\log r/2),&r=|X|\ge1,\\
 0,&r<1.
 \end{cases}
\end{equation}
On the exterior this is equivalently
$j_0=e^{tu^2}\Phi(u)\cos(au)/(8\pi r^3)$, with
$u=\frac12\log r$. It is a signed density with respect to ordinary
Euclidean volume $\dd X$, not with respect to $\dd r$. The spherical
Jacobian is
\begin{equation}\label{eq:densityjacobian}
 \frac{4\pi r^2\dd r}{8\pi r^3}
 =\frac{\dd r}{2r}=\dd u,\qquad
 \int_{\R^3}j_0(t,a;X)\dd X=H(t,a).
\end{equation}
The subscript zero means auxiliary time $s=0$; Newman time $t$ remains
free. The empty initial ball is part of the auxiliary domain, not a
second inverted theta contribution. It imposes no boundary at $r=1$
after evolution begins.

\begin{proposition}[Conservative theta heat evolution]\label{prop:conservative}
For $\varepsilon>0$ define
\begin{equation}\label{eq:conservativeheat}
\boxed{\partial_sj_s=\varepsilon\Delta j_s,\qquad
j_s=K_{\varepsilon s}*j_0,\qquad
K_{\varepsilon s}(X)=
 \frac{e^{-|X|^2/(4\varepsilon s)}}{(4\pi\varepsilon s)^{3/2}}.}
\end{equation}
This is the $L^1$-continuous mild solution, smooth for $s>0$.
It preserves $H$ and all finite spectral derivatives:
\begin{equation}\label{eq:conservedjets}
\int_{\mathbb R^3}\partial_a^k j_s\dd X=\partial_a^kH(t,a),
\qquad \partial_tj_s=-\partial_a^2j_s.
\end{equation}
The theta sum in \eqref{eq:fullthetaflow} passes termwise through this
evolution in $L^1$. Let $Q_{t,N}$ be the first $N$ summands of
$Q$ in \eqref{eq:Q}, multiplied by $e^{t\log^2r/4}$, and let
$j_s^{(N)}$ be the corresponding heat-evolved density. Then
\[
\|j_s-j_s^{(N)}\|_1\le
 \frac18\int_1^\infty r^{3/2}(Q_t-Q_{t,N})\dd r.
\]
\end{proposition}

\begin{proof}
The tails of Lemma~\ref{lem:tails} put the initial density and its finite
parameter derivatives in $L^1\cap L^2$, with all finite spatial moments.
The jump at $r=1$ is admissible for a mild initial datum.
The explicit Gaussian kernel gives the stated solution
\cite[Section~5.1]{HunterPDE}. Its integral is one, so Fubini proves
conservation and its $L^1$ norm gives the error bound.
Dominated parameter differentiation commutes with convolution.
Initially $\partial_tj_0=-\partial_a^2j_0$, proving the second identity.
Finally the positive theta summands have total $L^1$ majorant
$\frac18\int r^{3/2}Q_t\dd r=H(t,0)$, justifying termwise evolution.
\end{proof}

There is no artificial reflecting surface at $r=1$ in this example:
positive $s$ may place density inside the ball. The full-space integral
in \eqref{eq:conservedjets} includes it. This local differential equation
has Gaussian spreading and no finite propagation speed.
It starts from the full signed density, so it is nontrivial even though
the initial data and operator are rotationally symmetric.

\subsection{An exact colliding control}\label{sec:control}

For any prescribed $T_c>0$, consider the positive source
\begin{equation}\label{eq:control}
 \Phi_*(u)=(64u^4+6)e^{-(4+T_c)u^2},\qquad
 H_*(t,a)=\int_0^\infty e^{tu^2}\Phi_*(u)\cos(au)\dd u.
\end{equation}
Its heat integral converges for $t<T_c+4$. At $t=T_c$, Gaussian
integration gives
\begin{equation}\label{eq:controlevent}
 H_*(T_c,a)=\frac{\sqrt\pi}{16}(a^2/4-6)^2e^{-a^2/16},
 \qquad A_*=2\sqrt6.
\end{equation}
Thus $H_*(T_c,A_*)=\partial_aH_*(T_c,A_*)=0$, with an exactly
double physical zero. Its exterior density is obtained from
\eqref{eq:densityinitial} by replacing the canonical amplitude
$Q_t(r)/(32\pi\sqrt r)$ by
$e^{tu^2}\Phi_*(u)/(8\pi r^3)$, where $u=\frac12\log r$.
It has the finite moments and parameter derivatives used by the heat
and prepared feedback arguments on its stated convergence domain.
The fixed example below uses $T_c=1/4$.

The control lacks the complete canonical square source and its theta
tail, but those are not separating premises of the generic heat or
feedback theorem. No global finite Newman threshold is claimed for this
family. Other smooth controls retain their own convergence and moment
domains; atomic inputs require a distributional formulation outside
these source-function statements \cite{CollisionControls}.

\subsection{A quantified collision diagnostic}

Write
\[
M_{2,k}=\int_{\mathbb R^3}|X|^2\partial_a^k j_0\dd X,\qquad
A_{m,k}=\int_{\mathbb R^3}|X|^m|\partial_a^k j_0|\dd X.
\]
At the spatial origin the heat formula gives
\begin{align}
(4\pi\varepsilon s)^{3/2}\partial_a^k j_s(0)
 &=\partial_a^kH(t,a)-\frac{M_{2,k}}{4\varepsilon s}
                         +\mathcal R_{s,k},\label{eq:heatdiagnostic}\\
|\mathcal R_{s,k}|&\le\frac{A_{4,k}}{32\varepsilon^2s^2}.
\end{align}
Indeed $|e^{-x}-1+x|\le x^2/2$ for $x\ge0$.
The simpler estimate $|e^{-x}-1|\le x$ yields, for every $s>0$,
\begin{equation}\label{eq:heatdiagnosticbound}
\left|(4\pi\varepsilon s)^{3/2}\partial_a^k j_s(0)
                         -\partial_a^kH(t,a)\right|
\le\frac{A_{2,k}}{4\varepsilon s}.
\end{equation}
Thus at $H(T,a)=H_a(T,a)=0$, both central components $j_s$ and
$\partial_a j_s$ are $O(s^{-5/2})$; a nonzero conserved component has
leading order $s^{-3/2}$. A simple zero removes the leading term only
from the first component. This is a direction-independent diagnostic
of the two cancellations, with an explicit moment allowance.
It is a faster-decay statement at the specified observation, not a
growing instability or an all-height uniform estimate.

The exact control in Section~\ref{sec:control} has the same two cancellations and finite
moments on its actual convergence domain, so it has the same extra
decay. The other applicable smooth controls also pass this generic
conservation mechanism. No premise of the heat theorem distinguishes
theta from those controls. Details, the radial convolution formula and
the fixed-kernel qualification are in \cite{ConservativeNote}.

\section{An energy-normalized conservative feedback}\label{sec:feedback}

Ordinary heat makes a collision visible through additional decay.
The following construction instead forces a finite-time singularity
of a specified response. It is the fixed candidate of
\cite{FeedbackNote}; its global regularity is a pointwise noncollision
reformulation, not an independent theta estimate.

\subsection{Preparation, marked derivatives, and conservation}

Write $f(t,a;X)$ for the raw density in \eqref{eq:densityinitial} and
fix the preparatory smoothing time to one:
\[
 f^\delta(t,a)=K_1*f(t,a),\qquad
 K_\sigma(X)=(4\pi\sigma)^{-3/2}e^{-|X|^2/(4\sigma)}.
\]
The smoothing is needed because the original jump at $r=1$ need not
belong to $H^1(\R^3)$. The prepared theta density and each finite
parameter derivative are Schwartz functions. The same holds for the
Gaussian-polynomial control on its convergence domain. The full unit
ball belongs to the evolved domain.

Fix a mark $p=(T,A)$, where $A$ is the physical frequency. Construct a common family
$j_s^p(t,a;X)$ with initial value $f^\delta(t,a;X)$ and put
\begin{equation}\label{eq:feedbackpair}
 U_s=j_s^p(T,A),\qquad
 V_s=\left.\partial_a j_s^p(T,a)\right|_{a=A},\qquad W_s=(U_s,V_s).
\end{equation}
All unqualified spatial integrals in this and the next two sections
are over $\R^3$. Choose
\begin{align}
 v(X)&=X+\frac{X}{1+|X|^2},&
 d(X)&=\operatorname{div}v
       =3+\frac{3+|X|^2}{(1+|X|^2)^2},\label{eq:feedbackdrift}\\
 E(s)&=\|W_s\|_2^2,&G(s)&=\|\nabla W_s\|_2^2,\qquad
 D(s)=\int d|W_s|^2.
\end{align}
The equation and scalar feedback are
\begin{equation}\label{eq:feedbackpde}
 \boxed{\partial_sj=\Delta j+\alpha(s)\operatorname{div}(vj),
 \qquad \alpha(s)=\frac{2G(s)}{D(s)}.}
\end{equation}
Diffusivity, spatial scale, preparatory time, and relative weight of
the two components have all been fixed to one. The drift is radial,
but its profile is a design choice, not an identity of theta arithmetic.
The dependence of $\alpha$ on spatial integrals makes the feedback
nonlocal in the state.

First solve the pair $W$; its resulting scalar coefficient then defines
one linear propagator for the running parameters $(t,a)$, with $p$
held fixed. In particular, $\alpha$ is independent of the running
frequency $a$ in this family. Uniqueness
identifies the marked value and actual derivative with
\eqref{eq:feedbackpair}. At fixed $p$, that propagator preserves
\begin{equation}\label{eq:feedbackobservations}
 \int j_s^p(t,a)=H(t,a),\qquad
 \int\partial_a^kj_s^p(t,a)=\partial_a^kH(t,a),\qquad
 \partial_tj_s^p=-\partial_a^2j_s^p.
\end{equation}
Indeed the flux is $-\nabla j-\alpha vj$, so
\[
 \frac{d}{ds}\int_{|X|<R}j
   =\int_{|X|=R}(\partial_rj+\alpha v_rj)\,dS.
\]
Smoothness at the origin and weighted decay at infinity remove only
those limiting boundaries; flux through $r=1$ is retained. The same
argument applies to the finite parameter derivatives on their
convergence domains.

Moving the mark is a different differentiation. For
$z=\partial_Aj_s^p(t,a)$ with $(t,a)$ held fixed,
\begin{equation}\label{eq:feedbackmarkresponse}
 \partial_sz=\Delta z+\alpha\operatorname{div}(vz)
               +\alpha_A\operatorname{div}(vj),\qquad z_0=0,
 \qquad \alpha_A=\frac{2G_A}{D}-\frac{2GD_A}{D^2}.
\end{equation}
The diagonal frequency derivative is $V+z$, not $V$. The additional
response has zero complete integral. Different marks generally have
different feedback paths; one path is not claimed to magnify every
frequency simultaneously.

Direct differentiation gives, with $q=|X|^2$,
\begin{equation}\label{eq:feedbackcoefficients}
 \frac78I\le Dv\le2I,\qquad 3\le d\le6,\qquad
 \Delta d=\frac{2q^2+20q-30}{(1+q)^4},\qquad |\Delta d|\le30.
\end{equation}
Integration by parts therefore yields
\begin{equation}\label{eq:feedbackenergy}
 E'=-2G+\alpha D=0,\qquad E=E_0>0,\qquad
 3E_0\le D\le6E_0,\qquad
 \frac{G}{3E_0}\le\alpha\le\frac{2G}{3E_0}.
\end{equation}
In particular the denominator is positive even at a collision.

\subsection{Local construction and the exact gradient identity}

\begin{proposition}[Construction and continuation]\label{prop:feedbackexistence}
For nonzero Schwartz initial data $W_0$, the pair equation has a unique
local $H^1$ solution, smooth with the weighted decay needed above on
compact intervals of its maximal lifespan $[0,S_{\max})$.
Let $B(s)=\int_0^s\alpha(q)\dd q$. Then
\begin{equation}\label{eq:feedbackgradientbound}
 G(s)+3E_0\le(G(0)+3E_0)e^{10B(s)}.
\end{equation}
A finite maximal endpoint requires $B(s)\to\infty$; a bound for $G$
through a finite endpoint gives continuation.
\end{proposition}

\begin{proof}
Set $w(X)=X/(1+|X|^2)$ and remove only the affine part of the drift by
$W_s(X)=b^3Z_\tau(bX)$, $d\tau/ds=b^2$. For a nonzero $Z\in H^1$ set
\[
 e=\|Z\|_2^2,\quad g=\|\nabla Z\|_2^2,\quad
 b=(E_0/e)^{1/3},\quad
 \widetilde D=\int d(Y/b)|Z(Y)|^2\dd Y.
\]
The transformed equation is
\begin{equation}\label{eq:feedbacktransformed}
 \partial_\tau Z=\Delta Z+
 \frac{2bg}{\widetilde D}\operatorname{div}_Y
                    \{w(Y/b)Z\},\qquad Z_0=W_0.
\end{equation}
On $H^1$ balls with $e$ bounded away from zero, this nonlinearity is
locally Lipschitz from $H^1$ to $L^2$: $3e\le\widetilde D\le6e$,
and the multiplier, its spatial derivatives and its $b$ derivatives
are uniformly bounded while $b$ stays in a compact positive interval.
The heat bound $L^2\to H^1$ is at most $C(1+\tau^{-1/2})$ on bounded
time intervals. Its integrable singularity makes the Duhamel map
contractive on $C([0,h];H^1)$ for small $h$; see the heat-smoothing
argument in \cite[Lemmas 5.50--5.51]{HunterPDE}.

In these variables $e_\tau=-6ge/\widetilde D$, hence
$b_\tau/b=2g/\widetilde D$ and $b_s/b=2b^2g/\widetilde D=\alpha$.
Thus the transformation recovers exactly \eqref{eq:feedbackpde}.
After the affine dilation the common linear equation has bounded
smooth drift on each compact lifespan interval. Its smoothing and
weighted estimates propagate the stated data and parameter derivatives.
The local construction also depends smoothly on the mark on common
compact lifespan intervals, giving \eqref{eq:feedbackmarkresponse}.

For later use retain the full gradient identity
\begin{align}
 G'&=-2N+\frac{2G}{D}\mathcal K,
       &N&=\|\Delta W\|_2^2,\label{eq:feedbackgradientexact}\\
 \mathcal K&=2\sum_{\ell=1}^2\int
       \nabla W_\ell^{\,T}Dv\,\nabla W_\ell
       +\int d|\nabla W|^2-\int(\Delta d)|W|^2.\nonumber
\end{align}
The diffusion term is nonpositive and \eqref{eq:feedbackcoefficients}
gives $G'\le10\alpha G+30\alpha E_0$, proving
\eqref{eq:feedbackgradientbound}. A bounded $G$ and the fixed positive
$E_0$ give a uniform local continuation interval through the same
construction. Since $B$ is nondecreasing, a finite maximal endpoint
therefore forces $B\to\infty$.
\end{proof}

For the radial prepared theta and control fields,
Lemma~\ref{lem:radialdriftsquare} strengthens the conditional gradient
allowance; the preceding construction also applies to nonradial data.

\subsection{A finite collision deadline and its converse}

\begin{theorem}[Collision-driven divergence]\label{thm:feedbackdeadline}
For the prepared theta pair, or the admitted smooth control pair, put
\[
 m=\int W_0=(H(T,A),H_a(T,A)),\qquad
 \mathsf P_0=\int |X|\,|W_0(X)|\dd X.
\]
If $m=0$, then
\begin{equation}\label{eq:feedbackdeadline}
 S_{\max}\le S_*:=
 \frac{60}{7(5\pi^2)^{2/5}}\frac{\mathsf P_0^{4/5}}{E_0^{2/5}}
 <\infty.
\end{equation}
At this finite endpoint $B$, $\alpha$, $G$, and
$\|W_s\|_\infty$ diverge. Moreover, within the stated initial-data class,
\begin{equation}\label{eq:feedbackequivalence}
 S_{\max}=\infty\quad\Longleftrightarrow\quad m\ne0.
\end{equation}
No nondegeneracy of the physical collision is required.
\end{theorem}

\begin{proof}
On a compact lifespan interval use the positive mass-preserving
transition kernel of the drift $-\alpha(s)v$ and diffusion $\sqrt2$.
Two trajectories with the same Brownian increments satisfy
\[
 \frac{d}{ds}|X_s-Y_s|^2
  \le-2\varkappa\alpha(s)|X_s-Y_s|^2,\qquad \varkappa=7/8.
\]
Thus the backward propagator contracts Lipschitz constants by
$e^{-\varkappa B(s)}$. This is the synchronous-coupling argument for
convex confinement discussed in \cite{EberleCoupling}; the scalar
coefficient is continuous on the interval, and $v$ is globally
Lipschitz with linear growth.

For the unnormalized transform
$\widehat W(\xi)=\int e^{-i\xi\cdot X}W(X)\dd X$, subtract the
backward exponential's value at zero and use $m=0$ to obtain
\[
 |\widehat W_s(\xi)|\le |\xi|\mathsf P_0e^{-\varkappa B(s)}.
\]
Plancherel split at radius $R$ gives
\begin{equation}\label{eq:feedbackfouriersplit}
 E_0\le\frac{\mathsf P_0^2e^{-2\varkappa B}R^5}{10\pi^2}
                         +\frac{G}{R^2}.
\end{equation}
Choose $R^5=5\pi^2E_0e^{2\varkappa B}/\mathsf P_0^2$ and use
\eqref{eq:feedbackenergy}. This yields
\[
 B'=\alpha\ge c_0e^{\beta B},\qquad
 c_0=\frac{(5\pi^2)^{2/5}}6\frac{E_0^{2/5}}{\mathsf P_0^{4/5}},
 \qquad \beta=\frac7{10}.
\]
Hence $e^{-\beta B(s)}\le1-\beta c_0s$, proving
\eqref{eq:feedbackdeadline}. Proposition~\ref{prop:feedbackexistence}
forces $B\to\infty$ at the finite endpoint; the last bound and
\eqref{eq:feedbackenergy} then force $\alpha,G\to\infty$.

For the converse let $h$ solve the same scalar equation with
$h_0=|W_0|$. Positivity gives $|W_s|\le h_s$ and conserved mass
$M_0=\int h_0$. Its second moment $R_h=\int|X|^2h$ satisfies
\begin{equation}\label{eq:feedbackenvelopemoment}
 R_h(s)\le e^{-2B(s)}R_h(0)
    +6M_0\int_0^se^{-2(B(s)-B(q))}\dd q.
\end{equation}
At a finite maximal endpoint dominated convergence gives $R_h\to0$.
For each radius $r>0$,
\[
 |m|\le\|W_s\|_1
 \le |B_r|^{1/2}E_0^{1/2}+R_h(s)/r^2.
\]
First taking the endpoint limit and then $r\downarrow0$ proves $m=0$
and also $\|W_s\|_1\to0$. Since
$E_0\le\|W_s\|_\infty\|W_s\|_1$, the amplitude diverges.
Together with the first implication this proves
\eqref{eq:feedbackequivalence}.
\end{proof}

The theorem supplies a quantitative magnifier, not a blow-up profile,
a rate in the remaining auxiliary time, or a continuation past the
singularity. The exact control of Section~\ref{sec:control}, with
$T_c=1/4$, has the same finite breakdown at $(T_c,A_*)$.
All generic hypotheses of the magnifier hold for that control.
Other admitted smooth controls retain their original moment domains;
atomic inputs are not included in this source-function statement.
An excluding estimate must provide additional canonical information.

\subsection{The fixed generic initializer comparison is insufficient}

There is a useful scoped failure even before any flow simulation.
From \eqref{eq:feedbackenergy} and \eqref{eq:feedbackgradientbound},
\begin{equation}\label{eq:feedbacklocallife}
 G'\le\frac{20}{3E_0}G(G+3E_0),\qquad
 S_{\rm local}=\frac1{20}\log\left(1+\frac{3E_0}{G(0)}\right)
\end{equation}
is a guaranteed lower lifespan, with existence for $s<S_{\rm local}$.
To state the assigned initial bounds, write
$f_k=\partial_a^kf(T,A)$ and
\[
 C_{j,k}=\int |X|^j|f_k(X)|\dd X,\qquad
 F_{4,k}=\int e^{-|X|^2/16}f_k(X)\dd X,\qquad C=C_{0,0}+C_{0,1}.
\]
The specified Gaussian and moment bounds are
\begin{align}
 \overline P&=\sqrt{C_{0,0}C_{2,0}}+
       \sqrt{C_{0,1}C_{2,1}}+\frac4{\sqrt\pi}C,
       &\mathsf P_0&\le\overline P,\nonumber\\
 \overline G&=\tfrac34(8\pi)^{-3/2}C^2,
       &G(0)&\le\overline G,\label{eq:feedbackinitializer}\\
 \underline E&=\frac{(24\pi)^{3/2}}{(16\pi)^3}
                       (F_{4,0}^2+F_{4,1}^2),
       &E_0&\ge\underline E.\nonumber
\end{align}
For the first bound use the mean radius $4/\sqrt\pi$ of $K_1$ and
Cauchy--Schwarz on each raw first moment. For the second use Young's
inequality and $\|\nabla K_1\|_2^2=\tfrac34(8\pi)^{-3/2}$.
For the third evaluate $K_3*W_0=K_4*(f_0,f_1)$ at zero and apply
Cauchy--Schwarz with $\|K_3\|_2^2=(24\pi)^{-3/2}$.

Let $\eta=(F_{4,0}^2+F_{4,1}^2)/C^2$ and $p_0=\overline P/C$.
Then $0\le\eta\le1$ and $p_0\ge1+4/\sqrt\pi>3$, since the raw
support has $r\ge1$. Substitution into the lower lifespan and upper
collision deadline gives, when $\eta>0$,
\begin{equation}\label{eq:feedbackpaymentfailure}
 \frac{\underline S_{\rm local}}{\overline S_*}
 =\frac7{1200}\left(\frac{15\sqrt{6\pi}}{256}\right)^{2/5}
   \frac{\eta^{2/5}}{p_0^{4/5}}
   \log\left(1+\frac{3\sqrt3}{2}\eta\right)
 <\frac7{1200}<\frac1{100}.
\end{equation}
Here $\underline S_{\rm local}$ replaces $E_0,G(0)$ in
\eqref{eq:feedbacklocallife} by $\underline E,\overline G$, and
$\overline S_*$ replaces $\mathsf P_0,E_0$ in
\eqref{eq:feedbackdeadline} by $\overline P,\underline E$.
The bound follows because the parenthesized factor is less than one,
$p_0^{4/5}>2$, and the logarithm is less than two. If $\eta=0$ there is
no positive lower-energy payment. These particular initial bounds
cannot certify existence through the assigned collision deadline,
even with exact arithmetic. This is not a bound on the true lifespan
or an obstruction to sharper estimates.

\section{Positive-carrier comparison and quantitative regularity}
\label{sec:carrier}

The first estimate tested for the feedback uses a positive carrier
for the complete theta source \cite{CarrierNote}. It gives explicit
regularity once a physical first-jet margin is independently known,
but does not produce that margin.

\subsection{The carrier and its relative-density dissipation}

At the fixed mark put $q_T(u)=e^{Tu^2}\Phi(u)$ and, on $r\ge1$,
$\rho_T(r)=q_T(\log r/2)/(8\pi r^3)$. Define
\[
 W^{\rm raw}=\mathbf1_{r\ge1}\rho_T
             (\cos(Au),-u\sin(Au)),\qquad
 h^{\rm raw}=\mathbf1_{r\ge1}\rho_T(1+u^2),\qquad u=\tfrac12\log r,
\]
and prepare both by $K_1$ convolution. Let $h$ solve the same scalar
equation as $W$, with the existing feedback path $\alpha(s)$.
Since $|W^{\rm raw}|\le h^{\rm raw}$, positivity gives
\begin{equation}\label{eq:carrierdomination}
 h>0,\qquad |W|\le h,\qquad
 M:=\int h=P_0(T)+P_2(T),\qquad
 P_j(T)=\int_0^\infty u^jq_T(u)\dd u.
\end{equation}
Writing $R=\int|X|^2h$ and $\mu=|m|$, one has
\begin{equation}\label{eq:carriermoment}
 R_0=\int_0^\infty e^{4u}(1+u^2)q_T(u)\dd u+6M,
 \qquad R'\le6M-2\alpha R.
\end{equation}
The $6M$ in the preparation comes from the second moment of $K_1$.

Set $Z=W/h$ and
\[
 \mathcal F=\int h|Z|^2=\int\frac{|W|^2}{h},\qquad
 \mathcal I=\int h|\nabla Z|^2.
\]
Subtracting the equations for $W$ and $h$ gives the quadratic
general-relative-entropy identity \cite[Lemma 1.1]{MichelMischlerPerthame}:
\begin{align}
 \partial_sZ&=\Delta Z+(2\nabla\log h+\alpha v)\cdot\nabla Z,\nonumber\\
 \partial_s(h|Z|^2)&=\Delta(h|Z|^2)
      +\alpha\operatorname{div}(vh|Z|^2)-2h|\nabla Z|^2,\nonumber\\
 \mathcal F(s)+2\int_0^s\mathcal I(q)\dd q&=\mathcal F(0).
 \label{eq:carrierentropy}
\end{align}
To justify complete-space integration, use nonnegative radial cutoffs.
As $|Z|\le1$, the Laplacian cutoff error is at most
$CR_{\rm cut}^{-2}M$ and the drift error is at most
$C\sup\alpha\int_{|X|\ge R_{\rm cut}}h$, on a compact lifespan
interval. Both vanish. Monotone convergence for the nonnegative
dissipation and continuity at the prepared initial data give
\eqref{eq:carrierentropy}.

Cauchy--Schwarz in the preparatory convolution and complete integration
give
\begin{equation}\label{eq:carrierinitialentropy}
 \mathcal F(0)\le\int_0^\infty q_T(u)
       \frac{\cos^2(Au)+u^2\sin^2(Au)}{1+u^2}\dd u\le P_0(T).
\end{equation}
The exact decomposition at every later time is
\begin{equation}\label{eq:carrierdecomposition}
 \mathcal F(s)=\frac{\mu^2}{M}
       +\int h_s\left|Z_s-\frac mM\right|^2.
\end{equation}
Its constant floor is $\mu^2/M$. Neither positivity of $M$ nor the
upper bound \eqref{eq:carrierinitialentropy} supplies a positive
lower bound for $\mu$. A decreasing positive quantity need not retain
a positive uniform floor.

\subsection{A sharp mass--energy--moment inequality}

\begin{lemma}\label{lem:carriersharp}
If $W\in L^1\cap L^2$, $E=\|W\|_2^2>0$, and $h\ge|W|$ is
nonnegative with finite $R=\int|X|^2h$, then
\begin{equation}\label{eq:carriersharp}
 R\ge c_3\frac{\|W\|_1^{7/3}}{E^{2/3}}
   \ge c_3\frac{|\int W|^{7/3}}{E^{2/3}},\qquad
 c_3=\frac37\left(\frac{15}{14\pi}\right)^{2/3}.
\end{equation}
The constant is sharp in this general comparison class.
\end{lemma}
\begin{proof}
For $L>0$, split the mass using $1-|X|^2/L^2$, discard its negative
part, and use Cauchy--Schwarz:
\[
 \|W\|_1\le a_3E^{1/2}L^{3/2}+R/L^2,\qquad
 a_3^2=\int_{|Y|<1}(1-|Y|^2)^2\dd Y=\frac{32\pi}{105}.
\]
Maximizing $\|W\|_1L^2-a_3E^{1/2}L^{7/2}$ gives
$L^{3/2}=4\|W\|_1/(7a_3E^{1/2})$ and the stated constant.
Equality is attained by $W=eC(1-|X|^2/L^2)_+$, $h=|W|$, for a
constant unit vector $e$: its mass, energy and moment are respectively
$8\pi CL^3/15$, $32\pi C^2L^3/105$ and $8\pi CL^5/35$.
This extremizer is not asserted to be a theta trajectory.
\end{proof}

\begin{proposition}[Regularity from an independent first-jet margin]
\label{prop:carriermargin}
If $\mu>0$ is independently established, put
$r_\mu=c_3\mu^{7/3}E_0^{-2/3}$. Then for every $s\ge0$,
\begin{align}
 B(s)&\le\tfrac12\log(R_0/r_\mu)+3Ms/r_\mu,
       \label{eq:carrierBbound}\\
 G(s)+3E_0&\le(G(0)+3E_0)(R_0/r_\mu)^5
                     \exp(30Ms/r_\mu).\label{eq:carrierGbound}
\end{align}
The response obeys $\alpha\le2G/(3E_0)$.
\end{proposition}
\begin{proof}
Before the maximal endpoint, Lemma~\ref{lem:carriersharp} gives
$R\ge r_\mu$. Dividing \eqref{eq:carriermoment} by $R$ and integrating
yields \eqref{eq:carrierBbound}; use $R(s)\ge r_\mu$ at the upper
endpoint. Equation~\eqref{eq:feedbackgradientbound} then gives
\eqref{eq:carrierGbound}. These bounds are finite on each finite
interval, so Proposition~\ref{prop:feedbackexistence} extends the
solution to every $s\ge0$.
\end{proof}

The exponential rate in this sufficient estimate deteriorates as
$\mu^{-7/3}$ with the other budgets fixed. This is not a sharp rate for
the actual trajectory. With budgets
$\mu\ge\underline\mu>0$, $E_0\le\overline E$, $M\le\overline M$,
$R_0\le\overline R$ and $G(0)\le\overline G$, one can replace
$r_\mu$ by $r_*=c_3\underline\mu^{7/3}\overline E^{-2/3}$ and the
other constants by their upper bounds. At time zero the same lemma
also gives
\begin{equation}\label{eq:carrierlowerenergy}
 E_0\ge c_3^{3/2}\underline\mu^{7/2}\overline R^{-3/2}>0,
\end{equation}
which supplies a uniform denominator for the response bound.

At a finite breakdown endpoint, \eqref{eq:carriermoment} implies
$R(s)\to0$, as in \eqref{eq:feedbackenvelopemoment}. The stronger
mass inequality gives the quantitative statements
\begin{align}
 \mathcal F(s)&\le\|W_s\|_1
       \le c_3^{-3/7}E_0^{2/7}R(s)^{3/7}\longrightarrow0,
       \label{eq:carrierentropycollapse}\\
 \|W_s\|_\infty&\ge c_3^{3/7}E_0^{5/7}R(s)^{-3/7}
                      \longrightarrow\infty.\label{eq:carrieramplitude}
\end{align}
For the exact control at $(T_c,2\sqrt6)$,
$P_0=9\sqrt\pi/4$, $P_2=21\sqrt\pi/32$ and
$M=93\sqrt\pi/32>0$, while $\mu=0$. Its initial relative energy is
positive but tends to zero by \eqref{eq:carrierentropycollapse}.
Thus the displayed positive-carrier premises and dissipation alone
do not provide a floor through the collision deadline. The control
lacks the complete theta arithmetic and tail, but this comparison
has not used those features as a separating premise.

\subsection{Theta specialization and the matched direct comparison}

Fix finite $T_+,A_+\ge0$ with
\[
 \vartheta=\frac{A_+^4P_4(T_+)}{24P_0(T_+)}<1.
\]
The inherited fourth-moment bound in Appendix~\ref{sec:fourthmoment}
gives, on $0\le T\le T_+$ and $0\le A\le A_+$,
\begin{equation}\label{eq:carrierthetamargin}
 H(T,A)-\tfrac A2H_a(T,A)\ge P_0(0)(1-\vartheta),\qquad
 \mu\ge\underline\mu:=
 \frac{P_0(0)(1-\vartheta)}{\sqrt{1+A_+^2/4}}>0.
\end{equation}
At $A=0$ the first inequality follows directly from positivity.
The finite complete-source budgets
\begin{align*}
 \overline M&=P_0(T_+)+P_2(T_+),&
 \overline E&=(8\pi)^{-3/2}\overline M^2,\\
 \overline G&=\tfrac34(8\pi)^{-3/2}\overline M^2,&
 \overline R&=\int_0^\infty e^{4u}(1+u^2)q_{T_+}(u)\dd u+6\overline M
\end{align*}
follow from the Gaussian Young bounds and source monotonicity.
They supply the uniform version of
\eqref{eq:carrierBbound}--\eqref{eq:carrierlowerenergy}.
Every point in this domain was already separated by
\eqref{eq:carrierthetamargin}; the flow gives no larger domain or
independent exclusion. The exact control has
$A_*^4P_4/P_0=123>24$ at its event, so it fails this direct-margin
condition. The conclusion concerns this tested comparison, not every
possible estimate from canonical moments or signed source relations.

\section{The canonical source relation and its residual}
\label{sec:sourcerelation}

The second estimate tested for the same feedback uses the unchanged
complete source's logarithmic derivative \cite{SourceRelationNote}.
The raw relation is exact, but preparation and evolution do not
preserve it with zero residual. The following calculations identify
the terms that a valid coupled estimate must retain.

\subsection{The fixed operator and its admissible extension}

On $r\ge1$, write
\[
 \rho_T(r)=\frac{Q_T(r)}{32\pi\sqrt r}
          =\frac{q_T(\log r/2)}{8\pi r^3},\qquad
 \ell_T(r)=\log\rho_T(r),\qquad a_0(r)=r\log r.
\]
Here $Q_T$ is the complete sum in \eqref{eq:Q} multiplied by the
physical heat weight; all its square coefficients remain one.
For the raw frequency derivatives put
\begin{align*}
 F_0&=\mathbf1_{r\ge1}\rho_T\cos(Au),&
 F_1&=-\mathbf1_{r\ge1}u\rho_T\sin(Au),&
 F_2&=-\mathbf1_{r\ge1}u^2\rho_T\cos(Au).
\end{align*}
The exterior operator
\begin{equation}\label{eq:sourceoperator}
 \mathcal A_T=a_0(r)(\partial_r-\ell_T'(r))
\end{equation}
satisfies
\begin{equation}\label{eq:sourcerawrelation}
 \mathcal A_TF_0=AF_1,\qquad
 \mathcal A_TF_1=F_1+AF_2,
\end{equation}
because $a_0/(2r)=u$. The second relation is the actual frequency
derivative of the first and does not close on two components.

To specify a smooth whole-space operator, use canonical reflection
to extend $\rho_T$ positively to $r>0$, as justified below.
Choose a smooth radial cutoff $\eta=0$ on $r\le1/2$, $\eta=1$ on
$r\ge3/4$, and define
\begin{equation}\label{eq:sourceextension}
 b(X)=\eta(r)\log r\,X,\qquad
 c_T(r)=-\eta(r)r\log r\,\ell_T'(r),\qquad
 \mathcal A_T=b\cdot\nabla+c_T.
\end{equation}
For the explicit coefficient bound in
Proposition~\ref{prop:radialcompactstrict}, fix $\eta(r)=\chi(4r-2)$,
where $\chi=0$ for $t\le0$, $\chi=1$ for $t\ge1$, and
$\chi(t)=e^{-1/t}/(e^{-1/t}+e^{-1/(1-t)})$ for $0<t<1$.
This changes the comparison operator near the origin, not the source
or feedback equation. All cutoff derivatives are retained below.
The identities \eqref{eq:sourcerawrelation} also hold distributionally
across $r=1$: the jump of $F_0$ has a surface delta multiplied by
$a_0(1)=0$, while $F_1,F_2$ vanish on that sphere.

\subsection{Preparation creates a nonzero canonical defect}

Distributional integration by parts gives, for a raw density $f$,
\begin{align}
 [\mathcal A_T,K_\delta*]f(X)
 =\int\big\{&(b(X)-b(Y))\cdot\nabla_XK_\delta(X-Y)\nonumber\\
 &+[c_T(X)-c_T(Y)+\operatorname{div}b(Y)]
                      K_\delta(X-Y)\big\}f(Y)\dd Y.
 \label{eq:sourcepreparecommutator}
\end{align}
In particular, the divergence term has a plus sign and the raw
surface jump has already been accounted for.

With the mark $p$ fixed, let
\[
 J_{2,s}=\left.\partial_a^2j_s^p(T,a)\right|_{a=A},\qquad
 J_{2,0}=K_1*F_2,
 \qquad L_s=\Delta+\alpha(s)(v\cdot\nabla+d).
\]
Then $\partial_sJ_2=L_sJ_2$. It is not in general equal to $-u^2U_s$
after smoothing. Define the relation vector and its residual by
\begin{equation}\label{eq:sourceresidual}
 \mathcal T_s=(AV_s,V_s+AJ_{2,s}),\qquad
 \mathcal R_s=\mathcal A_TW_s-\mathcal T_s.
\end{equation}
At preparation and at the unit sphere, respectively,
\begin{align}
 \mathcal R_0&=([\mathcal A_T,K_1*]F_0,
               [\mathcal A_T,K_1*]F_1),\label{eq:sourceinitialresidual}\\
 \mathcal R_s(1)&=(-AV_s(1),-V_s(1)-AJ_{2,s}(1))
                 =-\mathcal T_s(1).\label{eq:sourcespheretrace}
\end{align}
The latter follows from $b=c_T=0$ at $r=1$, and is independent of
the chosen inner extension.

\begin{proposition}[Actual-theta preparation defect]
\label{prop:sourcepreparationdefect}
For a unit vector $X_1$ and $j=2,4$ put
\begin{align}
 L_j(T)&=\int_{|Y|\ge1}K_1(X_1-Y)\rho_T(|Y|)u(Y)^j\dd Y\nonumber\\
 &=\int_0^\infty q_T(u)u^j k(e^{2u})\dd u,
 \qquad
 k(r)=(4\pi)^{-3/2}e^{-(1+r^2)/4}\frac{\sinh(r/2)}{r/2}.
 \label{eq:sourcemoments}
\end{align}
These numbers are finite and strictly positive at every real $T$.
For
$0<A^2\le3L_2(T)/(2L_4(T))$,
\begin{equation}\label{eq:sourcedefectlower}
 \mathcal R_{U,0}(1)\ge\tfrac34A^2L_2(T)>0,
 \qquad \mathcal R_{V,0}(1)\ge A L_2(T)>0.
\end{equation}
The fixed zero-residual relation therefore cannot hold on an open
frequency interval after preparation.
\end{proposition}

\begin{proof}
Spherical averaging of the heat kernel gives
\eqref{eq:sourcemoments}. Equations~\eqref{eq:sourcespheretrace} and
the raw derivatives show that the two initial trace components are
the integrals of $A u\sin(Au)$ and
$u\sin(Au)+Au^2\cos(Au)$ against $K_1(X_1-Y)\rho_T(|Y|)\dd Y$.
Using $|\sin x-x|\le|x|^3/6$ and
$|\cos x-1|\le x^2/2$ gives
\begin{equation}\label{eq:sourcedefecterror}
 |\mathcal R_{U,0}(1)-A^2L_2|\le A^4L_4/6,\qquad
 |\mathcal R_{V,0}(1)-2AL_2|\le\tfrac23A^3L_4.
\end{equation}
The asserted interval yields \eqref{eq:sourcedefectlower}.
The first component is an entire even function of $A$, by the
complete theta tails, with nonzero quadratic coefficient $L_2$.
Its possible positive real zeros are isolated, proving the last
assertion.
\end{proof}

This argument does not determine the trace at a hypothetical collision
among those exceptional frequencies, and does not exclude an
event-conditioned estimate. It does refute the unchanged zero-residual
substitution directly on actual canonical data.

\subsection{The full residual equation and the sphere cancellation}

Because the same $L_s$ acts on $U,V,J_2$, one has
\begin{equation}\label{eq:sourceresidualevolution}
 \partial_s\mathcal R=L_s\mathcal R+[\mathcal A_T,L_s]W.
\end{equation}
With repeated spatial indices summed and $c=c_T$, the complete
commutator is
\begin{align}
 [\mathcal A_T,L_s]={}&-2(\partial_jb_i)\partial_{ij}
  -(\Delta b+2\nabla c)\cdot\nabla-\Delta c\nonumber\\
 &+\alpha\big\{(b\cdot\nabla v-v\cdot\nabla b)\cdot\nabla
                 +b\cdot\nabla d-v\cdot\nabla c\big\}.
 \label{eq:sourcefullcommutator}
\end{align}
The terms without a derivative to their right are multiplication
operators. All derivatives of the cutoff in
\eqref{eq:sourceextension} are included. There is no derivative of
$\alpha$: this is a spatial commutator with the mark fixed.
On the exterior,
\[
 \partial_jb_i=\log r\,\delta_{ij}+X_iX_j/r^2,
 \qquad \Delta b=3X/r^2,
\]
so the second-order part is
\begin{equation}\label{eq:sourceprincipalcommutator}
 -2\big((\log r)\Delta+\partial_{rr}\big).
\end{equation}
If $v=\gamma(r)X$, $\gamma=1+1/(1+r^2)$, the first-order Lie
coefficient is $(r\gamma'\log r-\gamma)X$.
Writing $P_{s,q}$ for the common propagator along the fixed mark,
\begin{equation}\label{eq:sourceresidualduhamel}
 \mathcal R_s=P_{s,0}\mathcal R_0+
       \int_0^sP_{s,q}[\mathcal A_T,L_q]W_q\dd q.
\end{equation}
The weighted smoothness justifies these identities on compact lifespan
intervals. They do not give a sign or a smallness bound for the forcing.

The theta pair $W$ remains radial by rotational invariance and
uniqueness. For $r>1$, the exact relation gives
$\partial_rW=\ell_T'W+(\mathcal T+\mathcal R)/(r\log r)$.
Consequently the actual gradient energy is
\begin{equation}\label{eq:sourcegradientrepresentation}
 \begin{split}
 G={}&4\pi\int_0^1r^2|\partial_rW|^2\dd r\\
 &+4\pi\int_1^\infty r^2
   \left|\ell_T'W+\frac{\mathcal T+\mathcal R}{r\log r}\right|^2\dd r.
 \end{split}
\end{equation}
The numerator $\mathcal T+\mathcal R$ vanishes at $r=1$ by
\eqref{eq:sourcespheretrace} and is $O(r-1)$ there. This cancellation
keeps the exterior integral finite. On the interval in
Proposition~\ref{prop:sourcepreparationdefect}, however,
$\mathcal T_{U,0}(1)=-\mathcal R_{U,0}(1)\ne0$.
Removing $\mathcal R$ would produce a false divergence proportional
to $\int_1^{1+\varepsilon}(r-1)^{-2}\dd r$, although $W_0$ is
Schwartz. Separately estimating the two singular quotients also loses
the cancellation required for finiteness. The interior integral and
the actual second frequency derivative must be retained as well.

\subsection{The uncut coefficient fails at the origin}

The theta inversion identity $\Theta(r)=r^{-1}\Theta(1/r)$
\cite{RadialNotes}, together with
\[
 \Delta_3[r^{-1}h(1/r)]=r^{-5}(\Delta_3h)(1/r),
\]
and the inversion-invariant heat weight, gives the complete reflection
law
\begin{equation}\label{eq:sourcereflection}
 Q_T(r)=r^{-5}Q_T(1/r),\qquad
 \rho_T(r)=r^{-6}\rho_T(1/r),\qquad r>0.
\end{equation}
The first square term of the unchanged series, with its differentiable
Gaussian tail, gives
\[
 \rho_T(r)=\frac\pi4r^{3/2}e^{-\pi r^2+(T/4)(\log r)^2}
 \left(1-\frac3{2\pi r^2}+O(e^{-3\pi r^2})\right)
 \quad(r\to\infty).
\]
Reflection therefore yields, as $r\downarrow0$,
\begin{align}
 \rho_T(r)&=\frac\pi4r^{-15/2}e^{-\pi/r^2+(T/4)(\log r)^2}
       \left(1-\frac{3r^2}{2\pi}+O(e^{-3\pi/r^2})\right),\nonumber\\
 \ell_T'(r)&=\frac{2\pi}{r^3}-\frac{15}{2r}
                   +\frac{T\log r}{2r}+O(r),\nonumber\\
 c_T^{\rm uncut}(r)&=-\frac{2\pi\log r}{r^2}
       +\frac{15}{2}\log r-\frac T2(\log r)^2+O(r^2\log r).
       \label{eq:sourceoriginsingularity}
\end{align}
For a smooth radial prepared field with $U_0(0)\ne0$,
\begin{equation}\label{eq:sourceoriginLtwo}
 \mathcal A_T^{\rm uncut}U_0(r)
   =-\frac{2\pi U_0(0)\log r}{r^2}+O((\log r)^2)
   \notin L^2(4\pi r^2\dd r)\quad\hbox{near zero}.
\end{equation}
This occurs for actual theta data. Define the positive finite moments
\[
 O_j(T)=(4\pi)^{-3/2}\int_0^\infty
          q_T(u)u^j e^{-e^{4u}/4}\dd u,\qquad j=0,2.
\]
Then $\cos x\ge1-x^2/2$ gives
\[
 U_0(0;A)\ge O_0-A^2O_2/2\ge O_0/2>0
       \quad\hbox{when }A^2\le O_0/O_2.
\]
Thus the uncut, unweighted $L^2$ operator-energy argument is not
admissible on this interval. The smooth extension
\eqref{eq:sourceextension} repairs the domain issue, while keeping
its transition terms and the actual inner-ball energy. Other weights,
subtractions or operator domains are outside this scoped failure.
The carrier-moment collapse alone does not specify an $L^2$
concentration profile at the origin.

\subsection{What the fixed source premise distinguishes}

For another positive smooth exterior amplitude $\rho_*$ and its own
raw pair, the operator frozen to theta satisfies
\begin{equation}\label{eq:sourcecomparison}
 \mathcal A_TF_{*,0}-AF_{*,1}
 =\mathbf1_{r\ge1}r\log r\,\rho_*
       \big((\log\rho_*)'-\ell_T'\big)\cos(Au).
\end{equation}
For each fixed $A\ge0$ this vanishes identically exactly when
$\rho_*=C\rho_T$ on $r>1$; the isolated cosine zeros do not change
the conclusion by continuity. Normalization fixes $C$ when required.
Thus the literal fixed coefficient identifies the complete source up
to scale. Constructing a separate logarithmic derivative for each
positive source is generic and gives no such distinction.

The Gaussian-polynomial control fails this fixed theta coefficient
because its tail differs. Weighted-square and finite-prime controls
change the full square coefficients; reflected, smooth-shape and
positive-primitive controls lack the full canonical square source;
Euler-factor controls change the source and conductor. Their own
logarithmic-source relations remain available, and the applicable
smooth controls still satisfy the generic feedback mechanism.
Atomic controls are outside this smooth-amplitude formulation.
These are the premise distinctions of the repository's control
register \cite{CollisionControls}, not a sign estimate for the
evolved canonical field.

The stopped variant is precisely the propagation or substitution of
this fixed relation with zero residual. The next section retains
\eqref{eq:sourceresidualevolution} and
\eqref{eq:sourcegradientrepresentation} and absorbs its second-order
forcing in one finite weighted energy. The stronger signed control
through the collision deadline remains open.

\section{A finite weighted residual estimate}\label{sec:weightedclosure}

The source residual can be retained in a finite coupled energy estimate
\cite{WeightedClosureNote}. Fix the explicit cutoff in
\eqref{eq:sourceextension}, the original feedback, and the physical marks
\begin{equation}\label{eq:weightedmarks}
 \mathscr P=[0,1/2]\times[1,6],\qquad p=(T,A).
\end{equation}
Keep the actual second running frequency derivative $J_2$ and the full
residual $\mathcal R$ of \eqref{eq:sourceresidual}. Set
\begin{equation}\label{eq:weightedenergy}
 \omega=(1+r^2)^{-2},\qquad Z_{\mathrm R}=\omega\mathcal R,
 \qquad Y=\|Z_{\mathrm R}\|_2^2,\qquad Q_2=\|J_2\|_2^2,
 \qquad \mathcal F_{\mathrm R}=G+Y+Q_2.
\end{equation}
The subscript distinguishes this energy and weighted residual from the
carrier quantities in Section~\ref{sec:carrier}. They are denoted
$\mathcal F$ and $Z$ in the source note. No new evolved field or feedback
coefficient is introduced. All spatial integrals remain over $\R^3$.

\subsection{Coefficient bounds and prepared data}

Write $[\mathcal A_T,L_s]=\mathscr C_0+\alpha\mathscr C_1$, where
$\mathscr C_0$ and $\mathscr C_1$ are respectively the full constant and
$\alpha$ coefficients in \eqref{eq:sourcefullcommutator}. In particular,
their multiplication terms and all derivatives of the cutoff remain.
Define
\begin{align}
 A_2&=2\sup\omega|Db|_{\mathrm F},&
 A_1&=\sup\omega|\Delta b+2\nabla c_T|,\nonumber\\
 A_0&=\sup\omega|\Delta c_T|,&
 B_1&=\sup\omega|b\cdot\nabla v-v\cdot\nabla b|,\nonumber\\
 B_0&=\sup\omega|b\cdot\nabla d-v\cdot\nabla c_T|.
 \label{eq:weightedcoefficients}
\end{align}
The matrix norm is Frobenius. Suprema are over all space and, for
coefficients depending on $T$, over $0\le T\le1/2$.

\begin{lemma}[Admissible coefficients and data]\label{lem:weighteddata}
All constants in \eqref{eq:weightedcoefficients} are finite. The prepared
energies are continuous on $\mathscr P$, with $E_0>0$, $Y(0)>0$, and
\begin{equation}\label{eq:weightedinitialcompactness}
 X_0:=\max_{p\in\mathscr P}
       \left(1+\frac{\mathcal F_{\mathrm R}(0;p)}{E_0(p)}\right)<\infty.
\end{equation}
These are spatial-energy statements; no physical first-jet margin is
assumed.
\end{lemma}
\begin{proof}
Put $\beta=3/(2\pi)$. The complete series gives, for $r\ge1$,
\begin{align*}
 \rho_T(r)&=\frac\pi4r^{3/2}
   e^{-\pi r^2+(T/4)(\log r)^2}P(r),\\
 P(r)&=1-\beta r^{-2}+\tau(r),\qquad
 \tau(r)=\sum_{n\ge2}(n^4-\beta n^2r^{-2})
                   e^{-\pi(n^2-1)r^2}.
\end{align*}
Every summand in $\tau$ is positive, so $P\ge1-\beta>0$.
Through order three, differentiated tails obey
$\tau^{(k)}=O_k(r^ke^{-3\pi r^2})$: after extracting the displayed
exponential, the remaining sums are bounded by polynomials in $n$
against $e^{-\pi(n^2-4)}$. The remainder is independent of $T$.
It follows uniformly on the stated time interval that
\[
 c_T=2\pi r^2\log r-\tfrac32\log r
       -\tfrac T2(\log r)^2+e(r),\qquad
 e^{(k)}=O(r^{-2-k}\log r)\quad(0\le k\le2).
\]
Consequently $c_T=O(r^2\log r)$, $\nabla c_T=O(r\log r)$ and
$\Delta c_T=O(\log r)$ at infinity. The other coefficients in
\eqref{eq:weightedcoefficients} grow at most as $r^2\log r$;
multiplication by $\omega=O(r^{-4})$ bounds them. On $[1/2,1]$,
the positive reflected source \eqref{eq:sourcereflection} is smooth
uniformly in $T$, so its logarithmic derivatives and all fixed cutoff
terms are bounded. Both $b$ and $c_T$ vanish near the origin.

All weighted $L^1$ moments of $F_0,F_1,F_2$ and their needed parameter
derivatives are uniformly finite. Differentiated Gaussian convolution,
using $(1+|X|)^m\le(1+|X-Y|)^m(1+|Y|)^m$, maps these data continuously
into Schwartz space. The coefficient bounds give continuity of the
prepared residual norms as well. Since $F_0$ is nonzero near $u=0$
and the heat multiplier $e^{-|\xi|^2}$ never vanishes, $U_0\ne0$.
Thus $E_0>0$, and compactness gives a positive minimum for $E_0$ and
\eqref{eq:weightedinitialcompactness}.

The preparation residual remains exactly
\eqref{eq:sourceinitialresidual}. On the open cutoff ball,
$\mathcal R_{U,0}=-AV_0$. For $A\in[1,6]$, the raw $F_1$ is nonzero
on a short exterior interval. Its Gaussian convolution is real analytic;
if it vanished on the whole ball, analytic continuation and heat
injectivity would force $F_1=0$. Hence $Y(0)>0$. This does not assert
nonzero trace at every frequency on the sphere $r=1$.
\end{proof}

\subsection{Signed production and absorption of second derivatives}

For $l=\log\omega$, the exact residual equation gives
\begin{align}
 \partial_sZ_{\mathrm R}={}&\Delta Z_{\mathrm R}
   +(\alpha v-2\nabla l)\cdot\nabla Z_{\mathrm R}\nonumber\\
 &+(\alpha d+|\nabla l|^2-\Delta l-\alpha v\cdot\nabla l)Z_{\mathrm R}
   +\omega(\mathscr C_0+\alpha\mathscr C_1)W.
 \label{eq:weightedconjugation}
\end{align}
Set $h_\omega=d-2v\cdot\nabla l$. Adding the energies for
$G,Z_{\mathrm R},J_2$ yields the full signed identity
\begin{align}
 \mathcal F_{\mathrm R}'={}&-2N-2\|\nabla Z_{\mathrm R}\|_2^2
   -2\|\nabla J_2\|_2^2
   +2\int|\nabla l|^2|Z_{\mathrm R}|^2
   +2\langle Z_{\mathrm R},\omega\mathscr C_0W\rangle\nonumber\\
 &+\alpha\left\{\mathcal K+\int h_\omega|Z_{\mathrm R}|^2
   +\int d|J_2|^2
   +2\langle Z_{\mathrm R},\omega\mathscr C_1W\rangle\right\}.
 \label{eq:weightedsignedproduction}
\end{align}
Here $N$ and $\mathcal K$ are those of
\eqref{eq:feedbackgradientexact}. In the residual energy, the advection
contributes $-\int(\alpha d-2\Delta l)|Z_{\mathrm R}|^2$, canceling
the Laplacians of $l$ from the multiplication term. With $q=r^2$,
\begin{equation}\label{eq:weightedweightconstants}
 2|\nabla l|^2=\frac{32q}{(1+q)^2}\le8,\qquad
 h_\omega=11+\frac{q-5}{(1+q)^2}\le\frac{265}{24}.
\end{equation}
The maxima occur at $q=1$ and $q=11$, respectively. All integrations
are justified on compact lifespan intervals by inherited Schwartz
regularity and polynomial coefficient growth, including at $s=0$
after preparation. Large-radius cutoff errors vanish. There is no
artificial boundary at $r=1$; \eqref{eq:sourcespheretrace} and the
combined numerator in \eqref{eq:sourcegradientrepresentation} remain.

\begin{theorem}[Finite weighted residual closure]\label{thm:weightedclosure}
For each mark in \eqref{eq:weightedmarks}, throughout its smooth lifespan,
\begin{align}
 \mathcal F_{\mathrm R}'+N+2\|\nabla Z_{\mathrm R}\|_2^2
       +2\|\nabla J_2\|_2^2
 \le{}&[A_1+\alpha(10+B_1)]G\nonumber\\
 &+[8+A_2^2+A_1+A_0
       +\alpha(265/24+B_1+B_0)]Y\nonumber\\
 &+6\alpha Q_2+[A_0+\alpha(30+B_0)]E_0.
 \label{eq:weightedexpandedbound}
\end{align}
In particular, define
\begin{align}
 \kappa_0&=8+A_2^2+A_1+A_0,\nonumber\\
 \kappa_1&=\max\{10+B_1,265/24+B_1+B_0,6,30+B_0\},\nonumber\\
 \mathscr D_{\mathrm R}&=N+2\|\nabla Z_{\mathrm R}\|_2^2
                              +2\|\nabla J_2\|_2^2.
 \label{eq:weightedconstants}
\end{align}
Then
\begin{equation}\label{eq:weightedclosure}
 \mathcal F_{\mathrm R}'+\mathscr D_{\mathrm R}
       \le(\kappa_0+\kappa_1\alpha)(\mathcal F_{\mathrm R}+E_0).
\end{equation}
The constants are defined finite source budgets, not acquired numerical
values, and contain no dependence on the unknown gradient or lifespan.
\end{theorem}
\begin{proof}
Whole-space integration by parts gives
$\|D^2W\|_2^2=\|\Delta W\|_2^2=N$. Frobenius Cauchy--Schwarz and
the triangle inequality therefore give
\begin{align*}
 \|\omega\mathscr C_0W\|_2&\le A_2\sqrt N+A_1\sqrt G+A_0\sqrt{E_0},\\
 \|\omega\mathscr C_1W\|_2&\le B_1\sqrt G+B_0\sqrt{E_0}.
\end{align*}
Use $2A_2\sqrt{YN}\le N+A_2^2Y$ to absorb one of the two units of
diffusion $N$ in \eqref{eq:weightedsignedproduction}. The remaining
payments are
$2A_1\sqrt{YG}\le A_1(Y+G)$,
$2A_0\sqrt{YE_0}\le A_0(Y+E_0)$ and the same inequalities with
$A_1,A_0$ replaced by $B_1,B_0$. Since $\alpha\ge0$,
$\mathcal K\le10G+30E_0$, $d\le6$, and
\eqref{eq:weightedweightconstants} hold, these give
\eqref{eq:weightedexpandedbound}. The coefficient maxima give
\eqref{eq:weightedclosure}. No differentiated residual equation or
third spectral jet is required.
\end{proof}

\subsection{Integrated dissipation, local existence and scope}

\begin{proposition}\label{prop:weightedconsequences}
Let $B(s)=\int_0^s\alpha(q)\dd q$ and
$\Xi(s)=\kappa_0s+\kappa_1B(s)$. For $s<S_{\max}$,
\begin{equation}\label{eq:weightedintegrated}
 e^{-\Xi(s)}(\mathcal F_{\mathrm R}(s)+E_0)
 +\int_0^se^{-\Xi(q)}\mathscr D_{\mathrm R}(q)\dd q
 \le\mathcal F_{\mathrm R}(0)+E_0.
\end{equation}
Put $x_0=1+\mathcal F_{\mathrm R}(0)/E_0$ and
$C=\kappa_0+2\kappa_1/3$. Then
\begin{equation}\label{eq:weightedlifespan}
 1+\frac{\mathcal F_{\mathrm R}(s)}{E_0}
 \le\frac{x_0}{1-Cx_0s}\quad(0\le s<(Cx_0)^{-1}),
 \qquad S_{\max}(p)\ge(Cx_0(p))^{-1}\ge(CX_0)^{-1}.
\end{equation}
For every fixed spatial ball,
\begin{equation}\label{eq:weightedlocalresidual}
 \|\mathcal R(s)\|_{L^2(B_{R_b})}^2\le(1+R_b^2)^4Y(s).
\end{equation}
\end{proposition}
\begin{proof}
Multiply \eqref{eq:weightedclosure} by $e^{-\Xi}$ and integrate.
For $x=1+\mathcal F_{\mathrm R}/E_0$, the feedback bound
$\alpha\le2G/(3E_0)\le2\mathcal F_{\mathrm R}/(3E_0)$ gives
$x'\le[\kappa_0+(2\kappa_1/3)(x-1)]x\le Cx^2$.
Integrating the inequality for $1/x$ proves the comparison first for
$s<\min\{S_{\max},(Cx_0)^{-1}\}$. If $S_{\max}<(Cx_0)^{-1}$,
it bounds $G$ up to that endpoint, contradicting
Proposition~\ref{prop:feedbackexistence}. Compactness from
Lemma~\ref{lem:weighteddata} gives the uniform open interval. The last
bound uses $\omega^{-2}\le(1+R_b^2)^4$ on the ball.
\end{proof}

This closes one weighted energy and its integrated dissipation. It
does not control the unweighted residual at infinity or every higher
derivative. The bound conditional on $B$ is
$\mathcal F_{\mathrm R}(s)+E_0\le(\mathcal F_{\mathrm R}(0)+E_0)e^{\Xi(s)}$;
no independent bound for $B$ through the collision deadline has been
supplied. For this original closure, $\kappa_1\ge30$, whereas
\eqref{eq:feedbackgradientbound} already controls $G+3E_0$ with exponent
$10B$. That closure alone does not sharpen the gradient allowance.
The radial square in Section~\ref{sec:radialgain} supplies the additional
quantitative improvement.

The exact control at $(1/4,2\sqrt6)$ lies in $\mathscr P$ and remains
admitted by the same theta operator and coefficient budgets. Its
prepared fields are Schwartz for $T<17/4$, but its residual must include
the raw source mismatch:
\begin{equation}\label{eq:weightedcontrolresidual}
 \mathcal R_{*,0}=[\mathcal A_T,K_1*]W_*^{\rm raw}
       +K_1*(\mathcal A_TW_*^{\rm raw}-\mathcal T_*^{\rm raw}).
\end{equation}
The second summand vanishes for theta and not for this control. Its
local interval is compatible with its finite collision deadline in
Theorem~\ref{thm:feedbackdeadline}. The complete source enters the
coefficients and initial data, but the Young payments discard the
mixed terms' signed correlations. No new physical first-jet margin,
collision-free rectangle, numerical cost advantage or threshold bound
follows. The other controls retain the premise and witness-domain
qualifications of \cite{CollisionControls,WeightedClosureNote}.


\subsection{A radial improvement of gradient production}
\label{sec:radialgain}

The complete prepared theta fields and the exact control remain radial
by rotational invariance and uniqueness. Retaining the spatial variation
in $\mathcal K$ gives a stronger gradient allowance without changing
$\mathcal F_{\mathrm R}$ or its source operator \cite{RadialDriftNote}.
The gain uses radiality and smooth decay; it does not use an arithmetic
sign distinguishing theta from the control.

\begin{lemma}[Exact radial drift square]\label{lem:radialdriftsquare}
For the smooth radial pair $W$ throughout the inherited lifespan, put
\begin{align}
 c_v(q)&=\frac{11+5q}{(1+q)^2},&
 p_v(q)&=\frac{q(203+59q+69q^2+21q^3)}{4(1+q)^4},\nonumber\\
 \mathscr H_{\rm rad}(W)&=
 \int c_v(r^2)\left|X\cdot\nabla W+\tfrac32W\right|^2
 +\int p_v(r^2)|W|^2\ge0.\label{eq:radialremainder}
\end{align}
Then
\begin{equation}\label{eq:radialdriftsquare}
 \mathcal K=10G+\frac{21}{4}E_0-\mathscr H_{\rm rad}(W).
\end{equation}
All integrals are over $\R^3$, including the inner ball.
\end{lemma}
\begin{proof}
Set $q=r^2$. The radial eigenvalue of $Dv$ is
$\nu=1+(1-q)/(1+q)^2$. Thus the gradient coefficient in $\mathcal K$
is
\[
 a_v=2\nu+d=5+\frac{5-q}{(1+q)^2},\qquad
 h_v=10-a_v=\frac{q(11+5q)}{(1+q)^2}.
\]
For a radial component $f$, completion of one square gives
\[
 \int_0^\infty r^2h_v|f'|^2\dd r
 =\int_0^\infty r^2h_v\left|f'+\frac{3f}{2r}\right|^2\dd r
 +\int_0^\infty r^2
  \left(\frac{3h_v'}{2r}-\frac{3h_v}{4r^2}\right)|f|^2\dd r.
\]
Here the prime on $h_v$ denotes its radial derivative. The boundary
term $rh_v|f|^2$ vanishes: $h_v=11r^2+O(r^4)$ at zero and the field
has the inherited decay at infinity. There is no boundary at $r=1$.
Using \eqref{eq:feedbackcoefficients}, the remaining potential is
\[
 -\Delta d-\frac{3h_v'}{2r}+\frac{3h_v}{4r^2}
 =\frac{21-119q+67q^2+15q^3}{4(1+q)^4}
 =\frac{21}{4}-p_v(q).
\]
Multiply by $4\pi$, sum over the components, and use
$h_v=r^2c_v(r^2)$ to obtain \eqref{eq:radialdriftsquare}.
The expression \eqref{eq:radialremainder} is smooth at the origin;
no singular operator or omitted interior contribution is involved.
\end{proof}

\begin{proposition}[Radial gradient and local compression bounds]
\label{prop:radiallocal}
For each mark in \eqref{eq:weightedmarks},
\begin{equation}\label{eq:radialgradientbound}
 G(s)\le\left(G(0)+\frac{21}{40}E_0\right)e^{10B(s)}
                  -\frac{21}{40}E_0.
\end{equation}
Put $g_0=G(0)/E_0>0$ and define
\begin{equation}\label{eq:radialclock}
 \delta(s)=1-\frac{4g_0}{3}(e^{7s/2}-1),\qquad
 L_{\rm rad}=\frac27\log\left(1+\frac{3}{4g_0}\right).
\end{equation}
Then $S_{\max}\ge L_{\rm rad}$, and on the open interval
$0\le s<L_{\rm rad}$,
\begin{equation}\label{eq:radiallocalbounds}
 \frac{G(s)}{E_0}\le\frac{g_0e^{7s/2}}{\delta(s)},\qquad
 \alpha(s)\le\frac{2g_0e^{7s/2}}{3\delta(s)},\qquad
 B(s)\le-\frac17\log\delta(s).
\end{equation}
The uniform lower lifespan follows by replacing $g_0$ in
\eqref{eq:radialclock} by
$\Gamma_0=\max_{p\in\mathscr P}G(0;p)/E_0(p)<\infty$.
\end{proposition}
\begin{proof}
The exact gradient identity and Lemma~\ref{lem:radialdriftsquare} give
\begin{equation}\label{eq:radialgradientproduction}
 G'+2N+\alpha\mathscr H_{\rm rad}
 =10\alpha\left(G+\frac{21}{40}E_0\right).
\end{equation}
Multiplication by $e^{-10B}$ and integration prove
\eqref{eq:radialgradientbound}. Next,
$G=-\langle W,\Delta W\rangle$ gives $N\ge G^2/E_0$.
With $g=G/E_0$ and $\alpha\le2g/3$, this yields
\[
 g'\le\frac{14}{3}g^2+\frac72g.
\]
Nonzero $L^2$ data have $g_0>0$. Solving this scalar comparison gives
the first two bounds in \eqref{eq:radiallocalbounds} before both the
lifespan and $L_{\rm rad}$. Integrating the alpha bound gives the last,
since $\delta'=-(14/3)g_0e^{7s/2}$. A finite maximal endpoint strictly
below $L_{\rm rad}$ would have bounded $G$, contradicting
Proposition~\ref{prop:feedbackexistence}. Lemma~\ref{lem:weighteddata}
gives continuity and compactness for the uniform assertion.
\end{proof}

At the same actual $B$, the allowance in
\eqref{eq:radialgradientbound} is smaller than
\eqref{eq:feedbackgradientbound} by
$(99/40)E_0(e^{10B}-1)$, strictly positive for $B>0$.
For a matched comparison, retaining $N\ge G^2/E_0$ under the old
$\mathcal K\le10G+30E_0$ gives the local interval
\[
 L_{\rm old,diff}=\frac1{20}\log\left(1+\frac{30}{7g_0}\right).
\]
For every $g_0>0$,
\begin{equation}\label{eq:radiallifespancomparison}
 L_{\rm rad}>L_{\rm old,diff}>
 \frac1{20}\log\left(1+\frac3{g_0}\right).
\end{equation}
Indeed, $\int_{g_0}^\infty(Ag^2+bg)^{-1}\dd g$ strictly increases
when $b$ decreases from $20$ to $7/2$, with $A=14/3$ fixed.
The final interval is the earlier diffusion-discarded comparison.
Direct analysis with the same $E_0,G(0)$ and radiality obtains this
same gain; residual coordinates acquire no additional first-jet
information or cost advantage.

\begin{proposition}[Refined radial weighted estimate]
\label{prop:radialweighted}
With precisely the previous functional, source budgets and Young
payments, the stronger expanded estimate is
\begin{align}
 \mathcal F_{\mathrm R}'+\mathscr D_{\mathrm R}
       +\alpha\mathscr H_{\rm rad}
 \le{}&[A_1+\alpha(10+B_1)]G\nonumber\\
 &+[\kappa_0+\alpha(265/24+B_1+B_0)]Y\nonumber\\
 &+6\alpha Q_2+[A_0+\alpha(21/4+B_0)]E_0.
 \label{eq:radialweightedexpanded}
\end{align}
Set
\[
 \kappa_{1,\rm rad}=\frac{265}{24}+B_1+B_0,\qquad
 \Xi_{\rm rad}(s)=\kappa_0s+\kappa_{1,\rm rad}B(s).
\]
Then
\begin{align}
 &e^{-\Xi_{\rm rad}(s)}(\mathcal F_{\mathrm R}(s)+E_0)
 +\int_0^s e^{-\Xi_{\rm rad}(q)}
    (\mathscr D_{\mathrm R}+\alpha\mathscr H_{\rm rad})(q)\dd q
 \le\mathcal F_{\mathrm R}(0)+E_0,\label{eq:radialweightedintegrated}\\
 &\mathcal F_{\mathrm R}(s)+E_0
 \le(\mathcal F_{\mathrm R}(0)+E_0)e^{\kappa_0s}
                   \delta(s)^{-\kappa_{1,\rm rad}/7}
 \quad(0\le s<L_{\rm rad}).\label{eq:radialweightedlocal}
\end{align}
The fixed-ball residual estimate \eqref{eq:weightedlocalresidual} remains.
\end{proposition}
\begin{proof}
Insert \eqref{eq:radialdriftsquare} in the full signed production
\eqref{eq:weightedsignedproduction} and repeat the unchanged absorption
proof of Theorem~\ref{thm:weightedclosure}. Only the $30+B_0$
coefficient of $\alpha E_0$ is replaced by $21/4+B_0$, giving
\eqref{eq:radialweightedexpanded}. The entry
$265/24+B_1+B_0$ dominates $10+B_1$, $6$ and $21/4+B_0$.
Multiplying the compact estimate by $e^{-\Xi_{\rm rad}}$ and
integrating proves \eqref{eq:radialweightedintegrated}; use
\eqref{eq:radiallocalbounds} for \eqref{eq:radialweightedlocal}.
\end{proof}

The expanded $\alpha E_0$ allowance decreases by exactly
$(99/4)\alpha E_0$. For the compact constant, the precise comparison is
\begin{equation}\label{eq:radialcompactgap}
 \kappa_1-\kappa_{1,\rm rad}
 =\max\{0,455/24-B_1\}.
\end{equation}
The fixed cutoff above pays this threshold analytically
\cite{RadialCompactNote}.

\begin{proposition}[Strict compact coefficient gain]
\label{prop:radialcompactstrict}
For the exact exponential cutoff $\eta(r)=\chi(4r-2)$ in
\eqref{eq:sourceextension}, with the unchanged drift and weight,
\begin{equation}\label{eq:radialcompactcoefficients}
\begin{aligned}
 B_1&\le\frac{69}{4},& \kappa_1&=30+B_0,\\
 \kappa_{1,\rm rad}&\le\frac{679}{24}+B_0,&
 \kappa_1-\kappa_{1,\rm rad}&\ge\frac{41}{24}.
\end{aligned}
\end{equation}
Define the old and radial compact upper allowances, for the same energy,
data and actual $B(s)$, by
\[
 U_j(s)=(\mathcal F_{\mathrm R}(0)+E_0)
          e^{\kappa_0s+\kappa_j B(s)},\qquad
 \kappa_{\rm old}=\kappa_1,\quad
 \kappa_{\rm rad}=\kappa_{1,\rm rad}.
\]
Then
\begin{equation}\label{eq:radialcompactfactor}
 \frac{U_{\rm rad}(s)}{U_{\rm old}(s)}
 =e^{-(455/24-B_1)B(s)}\le e^{-41B(s)/24},
\end{equation}
strictly below one if $B(s)>0$. In particular,
\begin{equation}\label{eq:radialcompactexplicit}
 \mathcal F_{\mathrm R}(s)+E_0
 \le(\mathcal F_{\mathrm R}(0)+E_0)
       e^{\kappa_0s+(679/24+B_0)B(s)}
\end{equation}
on the inherited lifespan. The coefficient acquisition is independent
of the physical mark and the source logarithmic derivative.
\end{proposition}
\begin{proof}
Write $v=\lambda X$, $\lambda=1+(1+r^2)^{-1}$, and
$b=\beta X$, $\beta=\eta\log r$. Direct differentiation gives
$(b\cdot\nabla)v-(v\cdot\nabla)b
=r(\beta\lambda'-\lambda\beta')X$. Hence
\[
 B_1=\sup_{r\ge0}p(r)
 \left|\eta\{r\lambda'\log r-\lambda\}
             -\lambda r\eta'\log r\right|,
 \qquad p(r)=\frac r{(1+r^2)^2}.
\]
The zero continuation at the origin is smooth. We use
$p\le1/2$, $\lambda\le2$, $r^2/(1+r^2)^2\le1/4$ and
$|r\lambda'|\le1/2$.

For $0<t\le1/2$, differentiating the displayed cutoff yields
\[
 0\le\chi'(t)\le e^{-1/t+1/(1-t)}
                 \{t^{-2}+(1-t)^{-2}\}\le8.
\]
Indeed, $e^{1/(1-t)}\le e^2$, and each term in
$e^{-1/t}\{t^{-2}+(1-t)^{-2}\}$ is at most $4e^{-2}$:
$x^2e^{-x}$ decreases for $x=1/t\ge2$, while
$(1-t)^{-2}\le4$. Reflection
$\chi(1-t)=1-\chi(t)$ gives the same bound on the other half.
Thus $0\le\eta'\le32$. The constant extensions are smooth and flat;
near $t=0$ the denominator tends to $e^{-1}>0$ and the numerator
$e^{-1/t}$ is flat, and the other endpoint follows by reflection.

For $r\le1/2$ the coefficient vanishes. On $1/2<r<3/4$,
$|\log r|\le\log2<1$, and the two contributions satisfy
\[
 p\eta\bigl(|r\lambda'||\log r|+\lambda\bigr)\le\frac54,
 \qquad
 p\lambda r|\eta'||\log r|
 =\lambda|\eta'||\log r|\frac{r^2}{(1+r^2)^2}\le16.
\]
Here $\log2<1$ follows by integrating $1/u<1$ from $1$ to $2$.
On $3/4\le r\le1$, $\eta=1$, $\eta'=0$ and
$|\log r|\le\log(4/3)<1$, giving $5/4$ again.
For every $r\ge1$, $\log r\le r$ implies $p\log r\le1/4$,
so the entire infinite tail is at most
$p\lambda+(p\log r)|r\lambda'|\le1+1/8=9/8$.
Smoothness includes both cutoff endpoints. Thus $B_1\le69/4$.

In \eqref{eq:radialcompactgap},
$455/24-B_1\ge455/24-69/4=41/24>0$.
The old maximum is therefore exactly $30+B_0$, and
$265/24+B_1+B_0\le679/24+B_0$.
Taking the ratio of the two allowances proves
\eqref{eq:radialcompactfactor}; $B\ge0$ and
\eqref{eq:radialweightedintegrated} give
\eqref{eq:radialcompactexplicit}.
\end{proof}

The ratio compares upper allowances, not actual solution energies.
All commutator,
preparation-residual and actual-jet terms, the cancellation at $r=1$,
and the evolved inner ball remain unchanged. The exact colliding
control satisfies the radial hypotheses, with both terms of
\eqref{eq:weightedcontrolresidual} retained. Its finite breakdown
is compatible with these longer local intervals. Neither
\eqref{eq:radialgradientbound} in the unknown $B$ nor
\eqref{eq:radialweightedlocal} before its singular endpoint controls
a collision deadline. No theta separator is supplied.



\subsection{Exterior source--residual coupling}\label{sec:sourcecoupling}

The full theta coefficient permits a further bounded improvement
\cite{SourceCouplingNote}. Keep the same physical rectangle, source,
preparation, cutoff, actual jets and energy. Write
$\Theta=\mathcal T=(AV,V+AJ_2)$ and use subscripts in/out only for
restrictions of volume integrals to $r<2$ and $r\ge2$.

\begin{proposition}[Source-coefficient compact gain]
\label{prop:sourcecoupling}
On $r\ge2$, the complete source has $c_T>0$ and
\begin{equation}\label{eq:sourceslope}
 2\le k:=\frac{v\cdot\nabla c_T-b\cdot\nabla d}{c_T}\le6.
\end{equation}
Its signed coupling to the actual residual satisfies
\begin{equation}\label{eq:sourceexteriorpayment}
 2\langle Z_{\mathrm R},\omega\mathscr C_1W\rangle_{\rm out}
 \le-2Y_{\rm out}+\frac{48}{625}G_{\rm out}
                   +\frac{12}{625}\|\Theta\|_{\rm out}^2.
\end{equation}
Define
\begin{equation}\label{eq:sourcecompactconstant}
 \kappa_{1,\rm src}=\frac{274}{25}+B_1+B_0
       =\kappa_{1,\rm rad}-\frac{49}{600},\qquad
 \Xi_{\rm src}(s)=\kappa_0s+\kappa_{1,\rm src}B(s).
\end{equation}
Then the unchanged whole-space quantities obey
\begin{equation}\label{eq:sourcecompactbound}
 \mathcal F_{\mathrm R}'+\mathscr D_{\mathrm R}
       +\alpha\mathscr H_{\rm rad}
 \le(\kappa_0+\kappa_{1,\rm src}\alpha)
                         (\mathcal F_{\mathrm R}+E_0).
\end{equation}
For $s<S_{\max}$ this gives
\begin{equation}\label{eq:sourceintegrated}
 \begin{split}
 &e^{-\Xi_{\rm src}(s)}(\mathcal F_{\mathrm R}(s)+E_0)\\
 &\quad+\int_0^s e^{-\Xi_{\rm src}(q)}
       (\mathscr D_{\mathrm R}+\alpha\mathscr H_{\rm rad})(q)\dd q
 \le\mathcal F_{\mathrm R}(0)+E_0.
 \end{split}
\end{equation}
The same-data upper allowance
$U_{\rm src}=(\mathcal F_{\mathrm R}(0)+E_0)e^{\Xi_{\rm src}}$
therefore satisfies
\begin{equation}\label{eq:sourceallowancecomparison}
 \frac{U_{\rm src}}{U_{\rm rad}}=e^{-49B/600},\qquad
 \kappa_{1,\rm src}\le\frac{2821}{100}+B_0,\qquad
 \kappa_1-\kappa_{1,\rm src}\ge\frac{179}{100}.
\end{equation}
The ratio is strict for $B>0$ and equals one at $B=0$.
\end{proposition}
\begin{proof}
First pay the complete differentiated tail. With $q=r^2\ge4$ and
$\beta=3/(2\pi)<1/2$, the exact factor in the proof of
Lemma~\ref{lem:weighteddata}, now written as a function of $q$, is
\[
 P(q)=1-\frac\beta q+\tau(q),\qquad
 \tau(q)=\sum_{n\ge2}\left(n^4-\frac{\beta n^2}{q}\right)
                       e^{-\pi(n^2-1)q}.
\]
For $x=\pi(n^2-1)q\ge36$, the absolute zeroth, $q$-weighted first
and $q^2$-weighted second derivatives of each summand are bounded by
$n^4e^{-x}$, $n^4(1+x)e^{-x}$ and
$n^4(x^2+2x+2)e^{-x}$, respectively. Each is at most
$2n^4e^{-x/2}$: the polynomial factors are at most $2x^2$, and
$e^{x/2}\ge x^4/384\ge x^2$.
Since $n^4\le16^{n-1}$ and $n^2-1\ge3+5(n-2)$,
\[
 2\sum_{n\ge2}n^4e^{-6(n^2-1)}
 \le\frac{32e^{-18}}{1-16e^{-30}}
 <64e^{-18}<\frac1{4096}<\varepsilon:=\frac1{100}.
\]
Thus $0\le\tau<\varepsilon$, $|q\tau_q|<\varepsilon$ and
$|q^2\tau_{qq}|<\varepsilon$ uniformly; these summable bounds also
justify both differentiations. Set $p_0=1-\beta/q$ and $t=\tau/p_0$.
Using $p_0\ge7/8$, $|qp_{0,q}|\le1/8$ and
$|q^2p_{0,qq}|\le1/4$, the quotient rule gives
\[
 |qt_q|\le\frac{64}{49}\varepsilon<\frac32\varepsilon,
 \qquad |q^2t_{qq}|\le\frac{632}{343}\varepsilon<2\varepsilon.
\]
For $z=2qP_q/P=z_0+2qt_q/(1+t)$, $z_0=2\beta/(q-\beta)$,
and $t\ge0$, it follows that
\[
 |z|\le\frac27+3\varepsilon<\frac13,\qquad
 |2qz_q|\le\frac{32}{49}+14\varepsilon+9\varepsilon^2<1.
\]
The last term retains the squared first derivative of $t$.

Put $L=-r\ell_T'=2\pi q-3/2-(T/2)\log r-z$. Uniformly for
$0\le T\le1/2$, $3<\pi<4$ and $\log r\le q$ give
\[
 L\ge6q-\frac{11}{6}-\frac q4\ge5q>0,\qquad
 rL'\le16q+1\le\frac{65}{4}q.
\]
Also $rL'-2L=3-T/2+T\log r+2z-2qz_q\ge13/12>0$.
Since $\log r\ge\log2\ge2/3$, these imply
\[
 c_T=(\log r)L>0,\qquad
 2\le\frac{rc_T'}{c_T}
      =\frac1{\log r}+\frac{rL'}L\le\frac{19}{4}.
\]
The lower bound for $\log2$ follows by pairing $1/x$ and $1/(3-x)$
on $[1,2]$. Here $\lambda\le6/5$ and
$-rd'=2q(q+5)/(1+q)^3\in[0,2]$, so
$k=\lambda rc_T'/c_T-rd'/L\in[2,29/5]\subset[2,6]$.
Division by $c_T$ has only been used on $r\ge2$.

On this exterior set $a_b=r\log r$ and
$a_1=r(r\lambda'\log r-\lambda)<0$. The full radial commutator is
$\mathscr C_1W=a_1W'-kc_TW$. The residual definition gives exactly
\begin{equation}\label{eq:sourcecompletion}
 \omega\mathscr C_1W=-kZ_{\mathrm R}
       +\omega\{(a_1+ka_b)W'-k\Theta\}.
\end{equation}
Because $|a_1|/a_b=\lambda/\log r-r\lambda'\le53/25<4\le2k$,
we have $|a_1+ka_b|\le ka_b$. Also $\log r\le r/2$ on $r\ge2$,
so $\omega a_b\le2/25$ and $\omega\le1/25$.
Writing $F_c$ for the final term in \eqref{eq:sourcecompletion}, use
$2Z_{\mathrm R}\cdot F_c\le k|Z_{\mathrm R}|^2+|F_c|^2/k$ and
$|F_c|^2/k\le2k\omega^2(a_b^2|W'|^2+|\Theta|^2)$.
This proves \eqref{eq:sourceexteriorpayment}, including the actual-jet
cost. In fact, since $A\le6$,
\[
 |\Theta|^2\le(A^2+2)|V|^2+2A^2|J_2|^2
             \le38|W|^2+72|J_2|^2.
\]

In the full signed production \eqref{eq:weightedsignedproduction},
retain all the original $\mathscr C_0$ payments and
$\mathscr H_{\rm rad}$ from \eqref{eq:radialdriftsquare}.
On $r<2$ keep the pointwise payment
$B_1|\nabla W|^2+(B_1+B_0)|Z_{\mathrm R}|^2+B_0|W|^2$.
The weight production there is at most $h_\omega(4)=274/25$;
outside it is at most $265/24$. The complete coefficients of the
restricted energies in the $\alpha$ allowance are
\[
\begin{array}{c|cc}
 &r<2&r\ge2\\ \hline
 G&10+B_1&10+48/625\\
 Y&274/25+B_1+B_0&217/24\\
 Q_2&6&6+864/625\\
 E&21/4+B_0&21/4+456/625
\end{array}
\]
Every exterior entry is less than $274/25$; all entries are at most
$\kappa_{1,\rm src}$. The unchanged non-$\alpha$ allowance is at
most $\kappa_0(\mathcal F_{\mathrm R}+E_0)$ after absorbing one
unit of $N$. This proves \eqref{eq:sourcecompactbound}.
The split partitions a previously derived volume integral: no indicator
is differentiated and no flux is added at $r=2$. The original cutoff,
whole inner ball and joint cancellation at $r=1$ remain.
Multiplication by $e^{-\Xi_{\rm src}}$ proves
\eqref{eq:sourceintegrated}. Finally
$265/24-274/25=49/600$, and $B_1\le69/4$ gives
$274/25+B_1\le2821/100$, proving
\eqref{eq:sourceallowancecomparison}.
\end{proof}

On each fixed ball the improved consequence is
\begin{equation}\label{eq:sourceballbound}
 \|\mathcal R(s)\|_{L^2(B_{R_b})}^2
 \le(1+R_b^2)^4
       \{(\mathcal F_{\mathrm R}(0)+E_0)e^{\Xi_{\rm src}(s)}-E_0\}.
\end{equation}
These are smaller upper allowances for the same actual $B$, not ratios
of actual energies. The independent gradient-only interval
$L_{\rm rad}$ is unchanged. Direct analysis with the same source
coefficients, fields and jets obtains the same improvement.

The theta coefficient is essential to \eqref{eq:sourceslope}, but the
resulting inequality holds for every admitted radial comparison field
with its true residual and common evolution. It never uses the
theta-only vanishing of the raw source mismatch. In particular the
exact control at $(1/4,2\sqrt6)$ remains admitted, with both terms of
\eqref{eq:weightedcontrolresidual}. Its finite deadline is unchanged.
The new gain therefore supplies no independent compression bound through
a deadline, physical first-jet floor or enlarged collision-free region.

\section{The remaining signed theta estimate}\label{sec:remaining}

The conservative magnifier is explicit, locally constructed, and
quantified by Theorem~\ref{thm:feedbackdeadline}. Its central missing
input remains an independent bound for $G$ or $\alpha$ through the
collision deadline, using the unchanged complete theta source.
Global regularity reformulates pointwise noncollision by
\eqref{eq:feedbackequivalence}; asserting it or checking known
noncolliding points does not supply that bound \cite{TargetClassification}.

The carrier estimate \eqref{eq:carrierGbound} needs an independently
supplied first-jet margin and recovers its already separated domain.
The fixed source relation acquires a nonzero preparation residual;
discarding it destroys the sphere cancellation, and the uncut operator
has an origin-domain failure. Theorem~\ref{thm:weightedclosure} now
resolves one narrower issue: its origin-admissible weighted energy
absorbs the second-order commutator forcing with the existing diffusion,
without a third spectral jet. The integrated and local consequences
are \eqref{eq:weightedintegrated}--\eqref{eq:weightedlocalresidual}.
The exact radial square further gives
\eqref{eq:radialgradientbound}--\eqref{eq:radialweightedlocal}, including
a strictly longer local interval than the old diffusion-retaining
comparison. Proposition~\ref{prop:radialcompactstrict} now acquires
$B_1\le69/4$ for the fixed exponential cutoff, making the compact
coefficient gain strict with gap at least $41/24$. These are generic
radial and cutoff estimates; they do not bound the compression integral
through the collision deadline.

Proposition~\ref{prop:sourcecoupling} uses the complete theta tail and
the signed source--residual product to lower the compact coefficient by
a further $49/600$, paying the actual second jet and whole inner region.
Its full signed starting point is \eqref{eq:weightedsignedproduction}.
This improves the residual allowance, but only uses coefficients of the
fixed theta operator. It remains valid for the exact colliding control
with its true residual; the theta-only vanishing of the raw source
mismatch has not been consumed. The next unresolved issue is how
information about the actual canonical trajectory or its prepared data
can strengthen the estimate beyond this control-admitting class.
Retaining $(\mathcal T+\mathcal R)(1)=0$ ensures domain compatibility;
that trace identity alone does not fix a whole-space production sign.

A further estimate may improve the joint production bound or use a
different admissible functional for the same flow. Each attempt must
specify its domain and quantitative comparison before calculation,
retain the actual spectral derivatives, complete source, preparation
residual, full commutator and whole evolving ball, and identify its
additional source input. Favorable cancellation need not occur in
each term separately. A useful improvement must change a stated
residual, gradient, compression or existence allowance, or reduce a
specified information requirement; a representation alone does not
supply the missing control. Neither all norm estimates nor all alternate
formulations are excluded by the present Young bound.

Any collision-excluding application must state why relevant controls
fail its complete premises on domains containing their witnesses
\cite{CollisionControls}, and compare with direct analysis using the
same source information. No sign for $ND-G\mathcal K$, independent
compression bound through the deadline, or stronger collision-free
domain is established here. The finite weighted closure does not imply
a closed hierarchy for arbitrary further derivatives.

An application to a hypothetical positive de Bruijn--Newman threshold
would additionally require the time, height, multiplicity and endpoint
coverage in the existing finite-event reduction \cite{ThresholdNote}.
No numerical evolution or new interval certificate is claimed.
Overview05 and Progress05 \cite{ConservativeProgram} own current routing;
the current continuation returns the still-missing source-specific
trajectory estimate to conceptual CTF-B02 work \cite{SourceContinuation}.
The coefficient proofs are complete \cite{RadialCompactNote,SourceCouplingNote}; the original
issued continuation \cite{RadialContinuation} and the older R4 brief
\cite{WeightedContinuation} retain their historical scopes.


\appendix
\section{The inherited fourth-moment restriction}\label{sec:fourthmoment}

Set
$P_j(t)=\int_0^\infty u^je^{tu^2}\Phi(u)\dd u$. The elementary
function
\[
G(\theta)=\cos\theta+\frac\theta2\sin\theta-1+\frac{\theta^4}{24}
\]
satisfies $G(0)=G'(0)=0$ and
$G''(\theta)=\theta(\theta-\sin\theta)/2>0$ for $\theta>0$.
Integration against the positive source proves
\begin{equation}\label{eq:moment}
H(t,a)-\frac a2H_a(t,a)>P_0(t)-\frac{a^4}{24}P_4(t),\qquad a>0.
\end{equation}
Thus every multiple real zero requires $a^4P_4(t)>24P_0(t)$.
Since
\[
\frac{\mathrm d}{\mathrm dt}\frac{P_4}{P_0}
=\frac{P_6P_0-P_4P_2}{P_0^2}>0,
\]
as follows by symmetrizing the positive double integral,
there is no simultaneous zero for
\[
0\le t\le T,\qquad
0<a\le\left(\frac{24P_0(T)}{P_4(T)}\right)^{1/4}.
\]
Note16 \cite{FluxNote} gives fully specified conservative bounds for
these moments. It also proves that the matched direct-$u$ argument has
identical constants. Improvement over the ordinary second-moment
radius is equivalent to $P_0P_4<6P_2^2$, not proved here. This
inherited bounded result is not a new outcome of the conservative flow.

\section*{Acknowledgments and verification status}

This manuscript was drafted for Edward Baker with substantial LLM
assistance. The retained conservative constructions were prepared with
OpenAI's GPT-6-astra in Codex, at the actual recorded reasoning effort
\emph{max}, on 27--28 October 2026. The inherited fourth-moment
comparison in Note16 was prepared with GPT-6.1-sol at effort
\emph{ultra} on 27 October 2026. The present editorial refocus, CTF-E05 request
R1, was carried out with GPT-6.1-sol at the actual recorded effort
\emph{ultra} on 28 October 2026. Edward Baker specified whole-integral
conservation as the organizing principle and proposed the program of
magnifying collisions without a prescribed spatial direction. The bounded
CTF-E06 cleanup on 29 October 2026 was carried out with GPT-6.1-sol,
actual recorded effort \emph{ultra}, chat0019; it clarifies normalization,
parameter roles, attribution, and routing without changing the proofs.
The weighted estimate was designed in Note39 with GPT-6-astra at effort
\emph{max} and proved in Note40 with GPT-6.1-sol at effort \emph{ultra},
including delegated same-model checks. Its CTF-E08 assessment and
integration were requested in the original design chat and carried out
with GPT-6-astra at actual recorded effort \emph{max} on 29 October 2026.
This is a different-model assessment of Sol's execution by the original
designer, with shared inherited premises, not independent specialist
review.

The proofs and internal assessments in Notes25,28--29,31 and
Reviews14,16--19 record the retained constructions and their scopes.
These are analytical results with the classical inputs identified below;
their LLM assessments do not constitute independent specialist review.
Note40 and Review22 record the weighted lemma and its integration scope.
Note42 adds the radial drift square: its conceptual choice and proof were
contributed by a delegated GPT-6-astra agent at recorded effort
\emph{max}; parent drafting/checking used GPT-6.1-sol at recorded effort
\emph{ultra}, chat0021. The present R4 assessment and manuscript
integration were carried out by GPT-6.1-sol, effort \emph{ultra}, in
that same chat on 29 October 2026, with a prior contributing Sol checker.
Review23 records this internal assessment with shared premises and
contributor overlap; it is not an independent specialist review.
Note44 proves the fixed cutoff coefficient bound and strict compact
gain with GPT-6.1-sol, actual effort \emph{ultra}, chat0022. Its
user-requested E08 assessment and integration in the same chat used
GPT-6.1-sol at effort \emph{ultra} with two read-only same-model
checks. The parent also prepared the proof; Review24 records this
contributor overlap and shared inherited premises, without independent
specialist review.
Note47's source--residual coupling was derived with GPT-6-astra at
actual effort \emph{max}, chat0023. The frozen R5 assessment and present
integration were completed with GPT-6-astra at actual effort \emph{max},
chat0024, on 29 October 2026. Review26 records a separate-chat same-model
internal assessment with shared inherited premises; no subagents,
cross-model assessment or specialist independence are claimed.
The build record records source checks and the built-in compiler's
actual result. Compilation does not validate a mathematical proof.
No numerical evolution of the nonlinear theta feedback or new interval
certificate is claimed.

\begin{thebibliography}{99}
\bibitem{ThetaCompanion}
E. Baker research archive, drafted with LLM assistance,
\emph{The Theta Transform in Three Dimensions}, Investigation05,
\path{book/main.tex}, Sections~1--2. Background companion; its scalar
theta--Mellin identity is classical (Riemann, 1859; DLMF~25.5.13--14).
Note35 records the normalization and attribution check.

\bibitem{ConservativeProgram}
E. Baker research archive, prepared with LLM assistance,
\emph{Research overview for conservative theta flow estimates} and
\emph{Research progress for conservative theta flow estimates},
Investigation05, Overview05 and Progress05, 29 October 2026.
The overview specifies the bounded contracts; Progress05 owns live status.

\bibitem{RadialNotes}
E. Baker research archive, prepared with LLM assistance,
\emph{Radial theta geometry and possible collision mechanisms}
and \emph{Radial flux and possible restrictions on collisions},
Investigation05, Notes01 and 14, 27 October 2026. Internal research notes.

\bibitem{FluxNote}
E. Baker research archive, prepared with LLM assistance,
\emph{Canonical flux bounds, a coupled moment restriction, and the direct comparison},
Investigation05, Note16, 27 October 2026. Internal research note;
assessment in Review09.

\bibitem{ConservativeNote}
E. Baker research archive, prepared with LLM assistance,
\emph{Total integral conservation and a direction independent heat flow},
Investigation05, Note25, 27 October 2026; internal assessment in Review14.

\bibitem{FeedbackNote}
E. Baker research archive, prepared with LLM assistance,
\emph{A conservative flow with collision-driven response blow-up},
Investigation05, Note28, 28 October 2026; same-assistant assessment
in Review16.

\bibitem{CarrierNote}
E. Baker research archive, prepared with LLM assistance,
\emph{Positive-carrier comparison and quantitative regularity},
Investigation05, Note29, 28 October 2026; same-assistant assessment
in Review17.

\bibitem{SourceRelationNote}
E. Baker research archive, prepared with LLM assistance,
\emph{The fixed canonical source relation and its preparatory commutator},
Investigation05, Note31, 28 October 2026; same-assistant assessment
in Review18. Integration in this manuscript is assessed in Review19.

\bibitem{WeightedClosureNote}
E. Baker research archive, prepared with LLM assistance,
\emph{A weighted residual and second-jet estimate} and
\emph{Verified weighted residual closure}, Investigation05,
Notes39--40, 29 October 2026. Note39: GPT-6-astra, effort max;
Note40: GPT-6.1-sol, effort ultra, internal same-model execution checks.
CTF-E08 assessment and integration in Review22 by the original designer;
no independent specialist review.

\bibitem{RadialDriftNote}
E. Baker research archive, prepared with LLM assistance,
\emph{CTF-B02 R4: one radial drift-production improvement},
Investigation05, Note42, 29 October 2026. Astra/max conceptual contribution;
Sol/ultra parent drafting and internal checks. Assessment/integration in
Review23, Sol/ultra; no independent specialist review.

\bibitem{RadialCompactNote}
E. Baker research archive, prepared with LLM assistance,
\emph{CTF-E07 R2: all-radius compact-coefficient bound},
Investigation05, Note44, 29 October 2026. GPT-6.1-sol, effort ultra;
user-requested E08 assessment and integration in Review24, same model
and parent contributor with shared premises, not specialist review.

\bibitem{SourceCouplingNote}
E. Baker research archive, prepared with LLM assistance,
\emph{CTF-B02 R5: exterior source/residual coupling}, Investigation05,
Note47, 29 October 2026. GPT-6-astra, effort max, chat0023;
same-model E08 R5 assessment and integration in Review26 and Note49,
chat0024. Internal LLM work with shared premises, not specialist review.

\bibitem{SourceContinuation}
E. Baker research archive, prepared with LLM assistance,
\emph{Continuation after source-coupling assessment and integration},
Investigation05, Note50, 29 October 2026. The next task is a bounded
conceptual choice about actual-trajectory information; no new candidate
is tested by the continuation.

\bibitem{CompactContinuation}
E. Baker research archive, prepared with LLM assistance,
\emph{Continuation after strict compact-gain integration},
Investigation05, Note46, 29 October 2026. Records the assessed result
and the next conceptual source-specific estimate; no further estimate
is executed by that continuation.

\bibitem{RadialContinuation}
E. Baker research archive, prepared with LLM assistance,
\emph{Continuation after radial drift integration}, Investigation05,
Note43, 29 October 2026. Fixed Sol analytical follow-through for the
existing coefficient $B_1$; not an executed additional estimate.

\bibitem{WeightedContinuation}
E. Baker research archive, prepared with LLM assistance,
\emph{Continuation after weighted residual closure}, Investigation05,
Note41, 29 October 2026. Integration record and bounded CTF-B02 R4
brief; the proposed next round is not an additional established estimate.

\bibitem{EberleCoupling}
A. Eberle, \emph{Reflection couplings and contraction rates for diffusions},
\href{https://arxiv.org/abs/1305.1233v3}{arXiv:1305.1233v3}, 2 June 2014.
Inspected introduction, pp.~1--2, on synchronous coupling for convex
confinement; reference annotation11 records the scope. The feedback
equation and collision deadline here are not imported from that paper.

\bibitem{MichelMischlerPerthame}
P. Michel, S. Mischler and B. Perthame,
\emph{General relative entropy inequality: an illustration on growth models},
J. Math. Pures Appl. \textbf{84} (2005), 1235--1260.
\href{https://www.ceremade.dauphine.fr/~mischler/articles/41Galentropy.pdf}
{Author manuscript}, 3 March 2005, Lemma~1.1 and equations (1.1)--(1.3).
Selected-passage provenance and application in reference annotation12.

\bibitem{CollisionControls}
E. Baker research archive, prepared with LLM assistance,
\emph{Collision controls and the screening they impose},
repository result summary12,
\path{result-summaries/12_collision_controls_and_screening.md}.
This register routes to the original control proofs and certificates;
each family's stated hypotheses and review scope are retained.

\bibitem{TargetClassification}
E. Baker research archive, prepared with LLM assistance,
\emph{Restating the target in new words and counting it as progress},
repository pattern01, \path{avoid/01_restating_the_target.md}.
Classification guidance, not a source of a new mathematical theorem.

\bibitem{HunterPDE}
J. K. Hunter, \emph{Notes on Partial Differential Equations},
revision 18 June 2014.
\href{https://www.math.ucdavis.edu/~hunter/pdes/pde_notes.pdf}
{Author's UC Davis notes}.
Inspected Theorem4.20, Section4.12, and the heat-kernel formulas
and Theorem5.5 in Section5.1; scope in reference annotation10.
Lemmas5.50--5.51 and the Plancherel convention in Theorem5.71
were additionally checked for the feedback construction; annotation11.

\bibitem{ThresholdNote}
E. Baker research archive, prepared with LLM assistance,
\emph{Rodgers--Tao with a positive-threshold contradiction hypothesis},
Investigation02, Note02, 26 October 2026. Internal exposition of the
previously recorded threshold reduction and its named external inputs;
not an upper-bound theorem from this manuscript.
\end{thebibliography}
\end{document}
